\chapter{Krzywizna: od transportu równoległego do tensorów}
\label{ch:krzywizna}
\index{krzywizna}

Co właściwie znaczy, że przestrzeń jest zakrzywiona?
Możemy wykonać pomiary \emph{wewnątrz} niej: poprowadzić
geodezyjne (odpowiedniki prostych), zmierzyć obwód małego
okręgu albo przenieść wektor wzdłuż zamkniętej drogi.
Na płaszczyźnie mały okrąg o promieniu $r$ ma obwód
$2\pi r$, a wektor po obejściu pętli wraca bez obrotu.
Na sferze obwód jest mniejszy, a na płaszczyźnie
hiperbolicznej większy; transport po pętli może obrócić
wektor. \emph{Krzywizna} opisuje właśnie te odchylenia
od geometrii euklidesowej. Bardzo małe otoczenie
każdego punktu wygląda w pierwszym przybliżeniu
płasko; krzywizna ujawnia się w kolejnych wyrazach
rachunku dla małych odległości.

Nie chodzi o samo wygięcie w większej przestrzeni.
Mały fragment walca w $\R^3$ można rozwinąć na płaszczyznę
bez rozciągania, więc jego wewnętrzna krzywizna jest
zerowa. Definicja przez metrykę i koneksję ma sens
także wtedy, gdy nie przedstawiamy rozmaitości jako
powierzchni zanurzonej w $\R^3$.

W rozdziale o koneksjach porównywaliśmy wektory w różnych
punktach za pomocą transportu równoległego. Pytanie,
czy wynik zależy od drogi, prowadzi do tensora Riemanna:
pełnego opisu krzywizny w punkcie. Krzywizna sekcyjna
bada jedną wybraną płaszczyznę styczną; na powierzchni
jest to krzywizna Gaussa. Tensor Ricciego sumuje
krzywizny płaszczyzn zawierających wybrany kierunek,
a krzywizna skalarna zbiera informację ze wszystkich
kierunków. Poniżej pokażemy, co te sumy mówią o
objętości małych obszarów.

Przez większość rozdziału $(M,g)$ jest rozmaitością
\emph{riemannowską}, a $\nabla$ jej koneksją Leviego-Civity
z Twierdzenia~\ref{thm:levi-civita}. Domyślnie rozmaitości
nie mają brzegu; powierzchnie z brzegiem i narożnikami
wprowadzimy osobno przy twierdzeniu Gaussa--Bonneta. W przykładach
Schwarzschilda i Kerra $g$ ma sygnaturę lorentzowską.
Te same wzory tensorowe nadal obowiązują, lecz wtedy
nie każda płaszczyzna ma określoną krzywiznę sekcyjną,
a ślad z metryką zawiera znaki ujemne.

\section{Co mierzy krzywizna? Tensor Riemanna}

\begin{definicja}[konwencja znaku]\label{def:tensor-krzywizny}
\index{konwencja znaku@konwencja znaku}
Niech $X,Y,Z$ będą polami wektorowymi. Definiujemy
\begin{equation}\label{eq:riemann-def}
 R(X,Y)Z=\nabla_X\nabla_YZ-\nabla_Y\nabla_XZ
                        -\nabla_{[X,Y]}Z,
 \qquad
 \mathcal R(X,Y,Z,W)=g(R(X,Y)Z,W).
\end{equation}
Symbol $R$ nazywamy \emph{operatorem krzywizny}, a
$\mathcal R$ \emph{tensorem Riemanna} o wszystkich
indeksach opuszczonych. Ta konwencja daje dodatnią
krzywiznę sekcyjną zwykłej sfery. W literaturze spotyka
się także przeciwny znak $R$; przy porównywaniu wzorów
trzeba zawsze sprawdzić definicję.
\end{definicja}

Pierwsze dwa składniki mierzą różnicę między różniczkowaniem
najpierw w kierunku $Y$, a potem $X$, oraz w przeciwnej
kolejności. Człon $\nabla_{[X,Y]}Z$ usuwa różnicę pochodzącą
tylko z tego, że same pola $X,Y$ nie muszą komutować.
W efekcie $R(X,Y)Z$ zależy \emph{wyłącznie} od wartości
$X,Y,Z$ w danym punkcie.

\begin{stwierdzenie}[tensorowość i współrzędne]
\label{prop:riemann-tensorowy}
Wyrażenie $R$ jest liniowe nad $C^\infty(M)$ w każdym
z trzech argumentów. Jeżeli
$\nabla_{\partial_i}\partial_j=\Gamma^a_{ij}\partial_a$,
to przy umowie
$R(\partial_i,\partial_j)\partial_b
 =R^a{}_{b i j}\partial_a$ mamy
\begin{equation}\label{eq:riemann-wspolrzedne}
 R^a{}_{b i j}
 =\partial_i\Gamma^a_{j b}-\partial_j\Gamma^a_{i b}
  +\Gamma^a_{i c}\Gamma^c_{j b}
  -\Gamma^a_{j c}\Gamma^c_{i b}.
\end{equation}
W szczególności tensor krzywizny zawiera pochodne
współczynników koneksji, a więc drugie pochodne metryki.
\end{stwierdzenie}

\begin{proof}
Dla dowolnej funkcji $f$ reguły koneksji i nawiasu dają
$\nabla_{fX}Y=f\nabla_XY$ oraz
$[fX,Y]=f[X,Y]-Y(f)X$. Zatem
\[
 \begin{split}
 R(fX,Y)Z
 &=f\nabla_X\nabla_YZ-\nabla_Y(f\nabla_XZ)
        -\nabla_{f[X,Y]-Y(f)X}Z\\
 &=fR(X,Y)Z-Y(f)\nabla_XZ+Y(f)\nabla_XZ
  =fR(X,Y)Z.
 \end{split}
\]
Antysymetria $R(X,Y)=-R(Y,X)$ daje to samo w drugim
argumencie. W trzecim argumencie rozwińmy
$\nabla_Y(fZ)=Y(f)Z+f\nabla_YZ$ i odpowiedni wyraz
z zamienionymi $X,Y$. Składniki zawierające pochodne
samego $Z$ znoszą się parami, natomiast pozostały
współczynnik przy $Z$ wynosi
$X(Yf)-Y(Xf)-[X,Y]f=0$. Dlatego $R(X,Y)(fZ)=fR(X,Y)Z$.

Ponieważ $[\partial_i,\partial_j]=0$, wstawiamy pola
współrzędnych do~\eqref{eq:riemann-def}:
\[
 \nabla_{\partial_i}(\Gamma^c_{jb}\partial_c)
 =(\partial_i\Gamma^a_{jb}
       +\Gamma^c_{jb}\Gamma^a_{ic})\partial_a.
\]
Po odjęciu wyrażenia z $i,j$ zamienionymi miejscami
otrzymujemy~\eqref{eq:riemann-wspolrzedne}.
\end{proof}

\begin{uwaga}[mały obwód i transport]\label{uwaga:holonomia}
Wybierzmy współrzędne normalne w $p$ i niezależne
$X_p,Y_p\in T_pM$. Obraz brzegu równoległoboku
$0,\varepsilon X_p,\varepsilon(X_p+Y_p),\varepsilon Y_p$
przez mapę normalną jest rzeczywistą zamkniętą pętlą.
Jej transport $P_\varepsilon:T_pM\to T_pM$ spełnia
\[
 P_\varepsilon v=v-\varepsilon^2R(X_p,Y_p)v+O(\varepsilon^3).
\]
Uzasadnijmy znak i rząd reszty. W trywializacji współrzędnych
piszemy $\Gamma_i=(\Gamma^a_{ib})_{a,b}$. Równanie
transportu ma postać $\dot V=-\Gamma_i(x)\dot x^iV$.
Ponieważ $\Gamma_i(0)=0$, na całej małej pętli
$V=v+O(\varepsilon^2)$, a więc
\[
 P_\varepsilon v-v
 =-\oint\Gamma_i(x)\,\dd x^i\,v+O(\varepsilon^4)
 =-\varepsilon^2(\partial_X\Gamma_Y-\partial_Y\Gamma_X)_p v
     +O(\varepsilon^3).
\]
Ostatnia równość wynika z liniowego rozwinięcia $\Gamma$;
składniki kwadratowe dają po całkowaniu $O(\varepsilon^3)$.
W punkcie $p$ iloczyny $\Gamma_i\Gamma_j$ znikają,
więc nawias jest macierzą $R(X_p,Y_p)$.
Odwrócenie obiegu zmienia znak głównego wyrazu.
Pole równoległoboku jest asymptotycznie równe
$\varepsilon^2|X_p\wedge Y_p|$; aby mówić o zmianie
\emph{na jednostkę pola}, dzielimy również przez
$|X_p\wedge Y_p|$ (dla pary ortonormalnej wynosi ono $1$).
Samo $\Gamma^a_{ij}\ne0$
nie świadczy o krzywiźnie: we współrzędnych biegunowych
płaskiej płaszczyzny symbole Christoffela są niezerowe,
ale wszystkie składniki $R$ znikają.
\end{uwaga}

\begin{figure}[htbp]
\centering
\begin{tikzpicture}[>=Latex,scale=1]
  \fill[manifoldblue!8] (-3.2,-.65)
    .. controls (-2.5,-1.08) and (-1.4,-.98) .. (-.58,-.35)
    .. controls (-.32,.02) and (-1.28,1.10) .. (-2.25,1.05)
    .. controls (-3.05,.9) and (-3.7,-.23) .. cycle;
  \draw[manifoldblue!65!black,thick] (-3.2,-.65)
    .. controls (-2.5,-1.08) and (-1.4,-.98) .. (-.58,-.35);
  \draw[manifoldblue!65!black,thick] (-.58,-.35)
    .. controls (-.32,.02) and (-1.28,1.10) .. (-2.25,1.05);
  \draw[manifoldblue!65!black,thick] (-2.25,1.05)
    .. controls (-3.05,.9) and (-3.7,-.23) .. (-3.2,-.65);
  \coordinate (A) at (-2.55,-.28);
  \coordinate (B) at (-1.55,-.38);
  \coordinate (C) at (-1.34,.40);
  \coordinate (D) at (-2.40,.52);
  \draw[manifoldgold!90!black,very thick,->]
    (A) .. controls (-2.20,-.38) .. (B);
  \draw[manifoldgold!90!black,very thick,->]
    (B) .. controls (-1.30,-.14) .. (C);
  \draw[manifoldgold!90!black,very thick,->]
    (C) .. controls (-1.78,.49) .. (D);
  \draw[manifoldgold!90!black,very thick,->]
    (D) .. controls (-2.62,.23) .. (A);
  \fill[manifoldred] (A) circle (2pt) node[below] {$p$};
  \draw[manifoldred,very thick,->] (A)--++(.08,.67)
      node[above] {$v$};
  \draw[manifoldgreen!75!black,very thick,->]
       (A)--++(.38,.55)
      node[right] {$P_\gamma v$};
  \node[manifoldblue!80!black,align=center] at (-2,-1.37)
      {zakrzywiona powierzchnia};
  \fill[manifoldblue!5] (1.1,-.85)--(3.55,-.85)
       --(4.05,.85)--(1.6,.85)--cycle;
  \draw[manifoldblue!60!black,thick]
       (1.1,-.85)--(3.55,-.85)--(4.05,.85)--(1.6,.85)--cycle;
  \draw[manifoldgold!90!black,very thick,->]
       (1.7,-.28)--(2.75,-.28)--(3.05,.42)
       --(2.0,.42)--(1.7,-.28);
  \fill[manifoldred] (1.7,-.28) circle (2pt) node[below] {$p$};
  \draw[manifoldred,very thick,->] (1.7,-.28)--++(0,.63)
       node[left] {$v=P_\gamma v$};
  \node[manifoldblue!80!black] at (2.6,-1.37)
      {płaska płaszczyzna};
\end{tikzpicture}
\caption{Transport wektora wzdłuż małej pętli.
Na zakrzywionej powierzchni wektor może wrócić obrócony;
na płaszczyźnie euklidesowej wraca w tej samej postaci.
Strzałki i pętla są schematyczne.}
\label{fig:krzywizna-holonomia}
\end{figure}

\section{Symetrie tensora i krzywizna sekcyjna}

\begin{twierdzenie}[symetrie i pierwsza tożsamość Bianchiego]
\label{thm:symetrie-riemanna}
Dla koneksji Leviego-Civity zachodzą
\begin{align}\label{eq:symetrie-riemanna}
 \mathcal R(X,Y,Z,W)&=-\mathcal R(Y,X,Z,W)
                     =-\mathcal R(X,Y,W,Z),\\
 \mathcal R(X,Y,Z,W)&=\mathcal R(Z,W,X,Y),\notag\\
 R(X,Y)Z+R(Y,Z)X+R(Z,X)Y&=0.\notag
\end{align}
Ostatnia równość to \emph{pierwsza tożsamość Bianchiego}.
\end{twierdzenie}

\begin{proof}
Antysymetria w $X,Y$ wynika wprost z definicji.
Pozostałe tożsamości są punktowe, więc wolno obliczyć je
w wygodnych współrzędnych. Wybierzmy współrzędne normalne
w punkcie $p$: $\Gamma^a_{ij}(p)=0$ i
$\partial_k g_{ij}(p)=0$. Znikanie pierwszych pochodnych
metryki wynika tu ze zgodności $\nabla g=0$ i $\Gamma(p)=0$.
Korzystając ze wzoru na symbole Christoffela
\eqref{eq:christoffel-metryka} oraz
\eqref{eq:riemann-wspolrzedne}, dostajemy w $p$
\begin{equation}\label{eq:riemann-drugie-pochodne}
 \begin{split}
 \mathcal R_{ij k l}
 :=\mathcal R(\partial_i,\partial_j,\partial_k,\partial_l)
 =\frac12\bigl(&\partial_i\partial_k g_{jl}
      -\partial_i\partial_l g_{jk}\\
      &-\partial_j\partial_k g_{il}
      +\partial_j\partial_l g_{ik}\bigr)(p).
 \end{split}
\end{equation}
Wyrazy kwadratowe w $\Gamma$ znikają w $p$.
Pierwszą i ostatnią parę indeksów można zamienić
miejscami, bo pochodne mieszane komutują i $g_{ij}=g_{ji}$.
Zamiana indeksów wewnątrz którejkolwiek pary zmienia
znak. Wreszcie po dodaniu wzoru
\eqref{eq:riemann-drugie-pochodne} dla trójek
$(i,j,k)$, $(j,k,i)$ i $(k,i,j)$ każdy składnik
drugich pochodnych znosi się z innym. To daje tożsamość
Bianchiego po sparowaniu z dowolnym $W$.
Tensorowość pozwala przenieść wynik z bazy
współrzędnych na dowolne wektory i dowolny punkt.
\end{proof}

\begin{definicja}[krzywizna sekcyjna]\label{def:krzywizna-sekcyjna}
\index{krzywizna sekcyjna@krzywizna sekcyjna}
Niech $\sigma\subset T_pM$ będzie dwuwymiarową
płaszczyzną z bazą $u,v$. W przypadku riemannowskim
\begin{equation}\label{eq:krzywizna-sekcyjna}
 K_p(\sigma)
 =\frac{g(R(u,v)v,u)}{|u|^2|v|^2-g(u,v)^2}
 =\frac{\mathcal R(u,v,v,u)}{|u\wedge v|^2}.
\end{equation}
Mianownik jest kwadratem pola równoległoboku
rozpiętego przez $u,v$. Dla ortonormalnej bazy
$u,v$ po prostu $K_p(\sigma)=g(R(u,v)v,u)$.
W sygnaturze lorentzowskiej ten sam iloraz określamy
tylko dla płaszczyzn \emph{niezdegenerowanych},
czyli gdy wyznacznik macierzy Grama w mianowniku
jest różny od zera.
\end{definicja}

\begin{proof}[Dlaczego wzór nie zależy od wybranej bazy]
Przy zmianie bazy płaszczyzny macierzą $A$ licznik
zmienia się przez czynnik $(\det A)^2$.
Wynika to z dwuliniowości i antysymetrii tensora
$\mathcal R$ osobno w dwóch parach argumentów:
każda para działa jak pole zorientowanego
równoległoboku. Mianownik jest wyznacznikiem
macierzy Grama, która przechodzi w $A^{\mathsf T}GA$,
więc także mnoży się przez $(\det A)^2$.
Iloraz jest zatem własnością samej płaszczyzny.
\end{proof}

\begin{uwaga}[intuicja]\label{uwaga:intuicja-sekcyjna}
Krzywizna sekcyjna bada, jak geometria zachowuje się
w wybranych \emph{dwóch} kierunkach stycznych. Na
sferze geodezyjne wychodzące z punktu zwykle zbliżają
się względem modelu płaskiego; w przestrzeni
hiperbolicznej rozchodzą się szybciej. W przestrzeni
o wymiarze $n\geq3$ różne płaszczyzny przez ten sam
punkt mogą mieć różne krzywizny. Krzywizna powierzchni
ma tylko jedną taką płaszczyznę, czyli całe $T_pM$.
\end{uwaga}

\section{Krzywizna Gaussa i \texorpdfstring{Theorema Egregium}{Theorema Egregium}}

Dla powierzchni $M^2\subset\R^3$ można najpierw pomyśleć
o tym, jak wygina się ona w przestrzeni. Nie jest to jednak
to samo co krzywizna wynikająca z pomiarów długości
\emph{na samej powierzchni}. Walec wygina się w $\R^3$,
ale można go rozwinąć na płaszczyźnie bez rozciągania.

Wyobraźmy sobie okrąg złożony z punktów odległych
o tę samą małą \emph{długość geodezyjną} $r$ od $p$.
W trzech modelach jego obwód jest odpowiednio
\[
 \begin{array}{c|c|c}
 K=R^{-2}\text{ na }S_R^2 & K=0\text{ na }\R^2
                         & K=-1\text{ na }\mathbb H^2\\ \hline
 2\pi R\sin(r/R) & 2\pi r & 2\pi\sinh r
 \end{array}
\]
(na sferze bierzemy $r<\pi R$). Pierwszy obwód jest
mniejszy od $2\pi r$, trzeci większy. Dla dowolnej
gładkiej powierzchni zachodzi przy $r\downarrow0$
\begin{equation}\label{eq:gauss-maly-okrag}
 L_p(r)=2\pi r\left(1-\frac{K(p)}6r^2+O(r^3)\right).
\end{equation}
Tak można \emph{zmierzyć} znak i wartość $K(p)$ bez
oglądania powierzchni z zewnątrz. Wzór wynika z
równania pól Jacobiego wzdłuż promieni; wyprowadzimy
jego odpowiednik dla objętości w
Sekcji~\ref{sec:ricci-krzywizna-skalarna} i uzasadnimy równanie
Jacobiego w Sekcji~\ref{sec:druga-wariacja-jacobi}.
Wymiar $K$ to odwrotność kwadratu długości: po
powiększeniu promienia dwa razy poprawka względna
w~\eqref{eq:gauss-maly-okrag} rośnie w przybliżeniu
czterokrotnie.

\begin{definicja}[druga forma podstawowa]
\index{druga forma podstawowa@druga forma podstawowa}
\label{def:druga-forma-podstawowa}
Niech $F:U\subset\R^2\to\R^3$ będzie regularną
parametryzacją powierzchni, $n$ lokalnie wybranym
jednostkowym wektorem normalnym, a $D$ zwykłą pochodną
w $\R^3$. Dla pól stycznych $X,Y$ definiujemy
\[
 \mathrm{II}(X,Y)=\langle D_XY,n\rangle,
 \qquad D_Xn=-S X.
\]
Operator $S$ jest operatorem kształtu. Z równości
$X\langle n,Y\rangle=0$ wynika
$\mathrm{II}(X,Y)=g(SX,Y)$; z braku torsji $D$
wynika symetria $\mathrm{II}$. Wartości własne $S$
to krzywizny główne $k_1,k_2$, a ich iloczyn
$K_{\mathrm G}=\det S=k_1k_2$ to \emph{krzywizna Gaussa}.
Zmiana $n$ na $-n$ zmienia znaki obu $k_i$,
ale nie zmienia ich iloczynu.
\end{definicja}

\begin{twierdzenie}[Gauss; Theorema Egregium]
\label{thm:egregium}
Dla powierzchni z metryką indukowaną z $\R^3$
krzywizna Gaussa równa się krzywiźnie sekcyjnej:
\begin{equation}\label{eq:egregium}
 K_{\mathrm G}(p)=K_p(T_pM)
 =\frac{\mathrm{II}_{11}\mathrm{II}_{22}
                  -\mathrm{II}_{12}^{2}}
             {g_{11}g_{22}-g_{12}^{2}}.
\end{equation}
W szczególności $K_{\mathrm G}$ można wyznaczyć
\emph{wyłącznie z metryki} $g$: lokalna izometria
powierzchni zachowuje krzywiznę Gaussa.
\end{twierdzenie}

\begin{proof}
Rozłóżmy pochodną otaczającą na część styczną
i normalną. Zgodnie ze Stwierdzeniem
\ref{prop:koneksja-projekcja} część styczna jest
koneksją Leviego-Civity metryki indukowanej, więc
\begin{equation}\label{eq:gauss-dekompozycja}
 D_XY=\nabla_XY+\mathrm{II}(X,Y)n,
 \qquad D_Xn=-S X.
\end{equation}
Pochodna $D$ w $\R^3$ ma zerową krzywiznę.
Policzmy część styczną równości
$D_XD_YZ-D_YD_XZ-D_{[X,Y]}Z=0$.
Pierwszy składnik na podstawie
\eqref{eq:gauss-dekompozycja} daje
\[
 (D_XD_YZ)^\top
 =\nabla_X\nabla_YZ-\mathrm{II}(Y,Z)SX.
\]
Występuje tu minus: różniczkując składnik
normalny $\mathrm{II}(Y,Z)n$, bierzemy jego część
styczną $\mathrm{II}(Y,Z)D_Xn$.
Po odjęciu analogicznego wyrazu z $X,Y$
zamienionymi otrzymujemy \emph{równanie Gaussa}
\begin{equation}\label{eq:gauss-rownanie}
 R(X,Y)Z=\mathrm{II}(Y,Z)SX
                       -\mathrm{II}(X,Z)SY.
\end{equation}
Wybierzmy ortonormalną bazę $e_1,e_2$ płaszczyzny
stycznej. Podstawiając $X=e_1$, $Y=Z=e_2$
i biorąc iloczyn skalarny z $e_1$, znajdujemy
\[
 g(R(e_1,e_2)e_2,e_1)
 =\mathrm{II}_{22}\mathrm{II}_{11}
      -\mathrm{II}_{12}\mathrm{II}_{21}=\det S.
\]
To jest pierwszy wzór~\eqref{eq:egregium}.
Dla bazy współrzędnych $F_1,F_2$ licznik
zamienia się w wyznacznik macierzy $\mathrm{II}$,
a mianownik jest wyznacznikiem macierzy $g$;
otrzymujemy drugi wzór. Wreszcie lewa strona
\eqref{eq:egregium} jest ilorazem z definicji
\ref{def:krzywizna-sekcyjna}, a tensor $R$
oblicza się z samej metryki przez
\eqref{eq:christoffel-metryka} i
\eqref{eq:riemann-wspolrzedne}. Jeśli $F$ jest
lokalną izometrią dwóch powierzchni, przenosi
metrykę, jej koneksję Leviego-Civity (jednoznaczność)
i $R$; przenosi zatem również $K$. To właśnie
niezwykły, \emph{wewnętrzny} charakter krzywizny
Gaussa stwierdza Theorema Egregium.
\end{proof}

\begin{przyklad}[płaszczyzna, walec, sfera i siodło]
\label{prz:egregium-cztery}
Na płaszczyźnie euklidesowej $K=0$. Walec o promieniu
$a>0$ ma lokalną parametryzację
\[
 F(u,v)=(a\cos(u/a),a\sin(u/a),v),
 \qquad F^*g_{\R^3}=\dd u^2+\dd v^2.
\]
Wobec twierdzenia $K=0$, choć walec jest zgięty
w przestrzeni otaczającej. Jego krzywizny główne
to $\pm1/a$ (zależnie od wyboru normalnej) i $0$.
Sfera o promieniu $a$ ma $S=\pm a^{-1}\mathrm{Id}$,
stąd $K=a^{-2}>0$. Dla wykresu
$z=(x^2-y^2)/2$ w punkcie $(0,0,0)$ pierwsza
forma podstawowa jest jednostkowa, a druga ma,
zależnie od normalnej, macierz
$\operatorname{diag}(1,-1)$. Dlatego $K(0)=-1$.
Znaki dodatni, zerowy i ujemny opisują więc
różne zachowanie \emph{metryki}, nie sam fakt
widocznego zgięcia w $\R^3$.
\end{przyklad}

\begin{figure}[htbp]
\centering
\begin{tikzpicture}[scale=1,>=Latex]
 \begin{scope}[shift={(-3.7,0)}]
  \shade[inner color=white,outer color=manifoldblue!43] (0,0) circle (1.02);
  \draw[manifoldblue!70!black,thick] (0,0) circle (1.02);
  \draw[manifoldblue!45] (-1.02,0)--(1.02,0);
  \draw[manifoldgold!85!black,very thick]
   plot[domain=15:165,samples=80] ({1.02*cos(35)*sin(\x)},{1.02*cos(\x)});
  \node at (0,-1.43) {$K>0$ sfera};
 \end{scope}
 \fill[manifoldblue!13] (-.93,-.78)--(.93,-.78)--(.93,.75)
  arc[start angle=0,end angle=180,x radius=.93,y radius=.27]--cycle;
 \draw[manifoldblue!75!black,thick]
  (-.93,-.78)--(-.93,.75) (.93,-.78)--(.93,.75);
 \draw[manifoldblue!75!black,thick] (0,.75) ellipse (.93 and .27);
 \draw[manifoldblue!75!black,thick] (-.93,-.78)
  arc[start angle=180,end angle=360,x radius=.93,y radius=.27];
 \draw[manifoldgold!85!black,very thick] (-.12,-.75)--(-.12,.74);
 \node at (0,-1.43) {$K=0$ walec};
 % Projection of (x,y,(x^2-y^2)/2) to (.7(x-y),.6(x+y)+.5z).
 \begin{scope}[shift={(3.6,0)}]
  \path[fill=manifoldblue!12,draw=manifoldblue!65!black]
   plot[domain=-1:1,samples=30] ({.7*(-1-\x)},{.6*(-1+\x)+.25*(1-(\x)^2)}) --
   plot[domain=-1:1,samples=30] ({.7*(\x-1)},{.6*(\x+1)+.25*((\x)^2-1)}) --
   plot[domain=1:-1,samples=30] ({.7*(1-\x)},{.6*(1+\x)+.25*(1-(\x)^2)}) --
   plot[domain=1:-1,samples=30] ({.7*(\x+1)},{.6*(\x-1)+.25*((\x)^2-1)}) -- cycle;
  \foreach \v in {-.75,-.5,-.25,0,.25,.5,.75} {
   \draw[manifoldblue!30,thin] plot[domain=-1:1,samples=30]
    ({.7*(\x-\v)},{.6*(\x+\v)+.25*((\x)^2-(\v)^2)});
   \draw[manifoldblue!30,thin] plot[domain=-1:1,samples=30]
    ({.7*(\v-\x)},{.6*(\v+\x)+.25*((\v)^2-(\x)^2)});
  }
  \draw[manifoldgold!85!black,very thick] plot[domain=-1:1,samples=50]
   ({.7*\x},{.6*\x+.25*(\x)^2});
  \draw[manifoldred!85!black,very thick] plot[domain=-1:1,samples=50]
   ({-.7*\x},{.6*\x-.25*(\x)^2});
  \fill[manifoldblue!85!black] (0,0) circle (1.6pt);
  \node at (0,-1.43) {$K<0$ siodło};
 \end{scope}
\end{tikzpicture}
\caption{Sfera ma dodatnią, walec zerową, a siodło ujemną krzywiznę Gaussa.
Złoty łuk na sferze jest rzutem łuku wielkiego okręgu. Na siodle
$z=(x^2-y^2)/2$ wyróżniono przekroje $y=0$ i $x=0$;
spotykają się w zaznaczonym punkcie, w którym $K=-1$.
Wartości $K$ wynikają z metryk tych powierzchni.}
\label{fig:krzywizna-znaki}
\end{figure}

\section{Ricci, krzywizna skalarna i przykłady modelowe}
\label{sec:ricci-krzywizna-skalarna}

\begin{definicja}[dwa ślady krzywizny]
\index{dwa slady krzywizny@dwa ślady krzywizny}
\label{def:ricci-skalarna}
\emph{Tensor Ricciego} jest śladem odwzorowania
$W\mapsto R(W,U)V$:
\begin{equation}\label{eq:ricci-def}
 \mathrm{Ric}(U,V)=\operatorname{tr}
              \bigl(W\longmapsto R(W,U)V\bigr),
 \qquad \mathrm{Ric}_{ij}=R^a{}_{j a i}.
\end{equation}
\emph{Krzywizna skalarna} to jeszcze jeden ślad,
tym razem względem metryki:
\begin{equation}\label{eq:scalar-def}
 \mathrm{Scal}=\operatorname{tr}_g\mathrm{Ric}
               =g^{ij}\mathrm{Ric}_{ij}.
\end{equation}
Nie mylmy $\mathrm{Scal}$ z operatorem $R$ ani
z wartością $K(\sigma)$ przypisaną płaszczyźnie.
\end{definicja}

Tensor Ricciego jest symetryczny. Istotnie, z pierwszej
tożsamości Bianchiego
$R(e_i,U)V-R(e_i,V)U=-R(U,V)e_i$.
Bierzemy iloczyn z $e_i$ i sumujemy po bazie ortonormalnej:
\[
 \mathrm{Ric}(U,V)-\mathrm{Ric}(V,U)
 =-\sum_i g(R(U,V)e_i,e_i)=0.
\]
Ostatnie składniki znikają przez antysymetrię w dwóch
ostatnich argumentach $\mathcal R$. W sygnaturze
pseudoriemannowskiej ta sama suma ma wagi
$\varepsilon_i=g(e_i,e_i)=\pm1$ i również znika.

Można to zobaczyć bez indeksów. Z punktu $p$ wysyłamy
geodezyjne w kierunku jednostkowego wektora $v$ i
w sąsiednich kierunkach. $\mathrm{Ric}(v,v)$ opisuje
\emph{łączny} wpływ krzywizn płaszczyzn rozpiętych
przez $v$ i każdy kierunek prostopadły na to, jak
rozszerza się wąska wiązka tych geodezyjnych.
Krzywizna skalarna zbiera te wyniki po wszystkich
kierunkach $v$. Ścisłe znaczenie słów ,,łącznie''
i ,,po wszystkich'' podają poniższe sumy.

\begin{stwierdzenie}[sumowanie krzywizn sekcyjnych]
\label{prop:ricci-sekcje}
Jeżeli $e_1,\ldots,e_n$ jest ortonormalną bazą
riemannowskiego $T_pM$, to
\begin{equation}\label{eq:ricci-suma-sekcji}
 \mathrm{Ric}(v,v)
 =\sum_{i=1}^n g(R(e_i,v)v,e_i).
\end{equation}
Dla jednostkowego $v=e_1$ jest to suma
$\sum_{i=2}^n K(\operatorname{span}\{e_1,e_i\})$.
Ponadto
\begin{equation}\label{eq:scalar-suma-sekcji}
 \mathrm{Scal}=2\sum_{1\leq i<j\leq n}
                    K(\operatorname{span}\{e_i,e_j\}).
\end{equation}
W wymiarze dwa: $\mathrm{Ric}=Kg$ oraz
$\mathrm{Scal}=2K$.
\end{stwierdzenie}

\begin{proof}
Ślad endomorfizmu w bazie ortonormalnej to suma
iloczynów skalarnych jego wartości na $e_i$ z $e_i$.
Stosując to do definicji~\eqref{eq:ricci-def},
dostajemy~\eqref{eq:ricci-suma-sekcji}. Składnik
z $e_i=v$ znika przez antysymetrię $R(v,v)=0$;
pozostałe są wartościami~\eqref{eq:krzywizna-sekcyjna}.
Ponowny ślad daje sumę po wszystkich parach
uporządkowanych $i\ne j$. Symetria par z
Twierdzenia~\ref{thm:symetrie-riemanna} sprawia,
że pary $(i,j)$ i $(j,i)$ mają tę samą krzywiznę;
stąd czynnik $2$. W wymiarze dwa dla jednostkowego
$v$ z pierwszej równości mamy
$\mathrm{Ric}(v,v)=K(p)$; jednorodność i polaryzacja
dają $\mathrm{Ric}=Kg$.
\end{proof}

\subsection*{Co mierzy Ricci, a co mierzy skalar?}
Załóżmy $n\geq2$. Ustalmy jednostkowy $v\in T_pM$ i patrzmy na geodezyjną
$\gamma_v(t)=\exp_p(tv)$. Sąsiednie promienie odchylone
w kierunkach prostopadłych do $v$ opisuje macierz pól
Jacobiego $A(t)$ z warunkami
$A(0)=0$, $A'(0)=\mathrm{Id}$; szczegóły konstrukcji
podamy w Sekcji~\ref{sec:druga-wariacja-jacobi}.
Ich równanie ma postać
$A''+\mathsf K_v(t)A=0$, gdzie
$\mathsf K_v(0)w=R(w,v)v$. Wstawiając warunki
początkowe i różniczkując równanie w zerze, otrzymujemy
\[
 A(t)=t\,\mathrm{Id}-\frac{t^3}{6}\mathsf K_v(0)+O(t^4).
\]
Dlaczego wyraz przy $t^2$ znika? Z równania
$A''(0)=-\mathsf K_v(0)A(0)=0$; trzecia pochodna
wynosi $A'''(0)=-\mathsf K_v(0)A'(0)=-\mathsf K_v(0)$.
Ślad $\mathsf K_v(0)$ to $\mathrm{Ric}(v,v)$, zatem
wyznacznik, czyli gęstość objętości w kierunku $v$,
spełnia
\begin{equation}\label{eq:ricci-gestosc-malej-kuli}
 j(t,v)=t^{n-1}\left(1-\frac{\mathrm{Ric}_p(v,v)}6t^2
                              +O(t^3)\right).
\end{equation}
Użyliśmy tu elementarnego rozwinięcia
$\det(\mathrm{Id}+\varepsilon B)
=1+\varepsilon\operatorname{tr}B+O(\varepsilon^2)$.
Przy dodatnim $\mathrm{Ric}(v,v)$ mała porcja objętości
wzdłuż tego promienia jest mniejsza niż w $\R^n$;
przy ujemnym jest większa. Mówimy o \emph{sumie}
odkształceń prostopadłych, nie o zachowaniu każdej
pojedynczej geodezyjnej.

Reszty są jednostajne względem $v\in S_pM$: zależność
rozwiązań od kierunku jest gładka, a sfera kierunków
jest zwarta. Dla dostatecznie małego $r$ kula metryczna
jest obrazem kuli stycznej przez $\exp_p$, zgodnie
z~\eqref{eq:male-kule-normalne}. Można zatem całkować
w tych współrzędnych bez wielokrotnego liczenia punktów.
Zsumujmy~\eqref{eq:ricci-gestosc-malej-kuli} po
wszystkich jednostkowych kierunkach. Symetria sfery
jednostkowej daje średnią
$\int_{S_pM}\mathrm{Ric}(v,v)\,\dd\theta/
\operatorname{area}(S^{n-1})=\mathrm{Scal}(p)/n$:
średnia $v_iv_j$ wynosi $\delta_{ij}/n$, a ślad
macierzy Ricciego to właśnie $\mathrm{Scal}$.
Całkując jeszcze promień $t$ od $0$ do $r$, przy
$b_n=\operatorname{vol}_{\R^n}(B_1(0))$ dostajemy
\begin{equation}\label{eq:scalar-mala-kula}
 \operatorname{vol}_g(B_g(p,r))
 =b_nr^n\left(1-\frac{\mathrm{Scal}(p)}{6(n+2)}r^2
                                      +O(r^3)\right).
\end{equation}
Faktycznie $\operatorname{area}(S^{n-1})=nb_n$,
a iloraz całek $\int_0^r t^{n+1}\dd t$ i
$\int_0^r t^{n-1}\dd t$ to $nr^2/(n+2)$.
Dodatnia krzywizna skalarna oznacza więc mniejszą
objętość bardzo małych kul niż w geometrii euklidesowej.
W wymiarze dwa $\mathrm{Scal}=2K$ i
\eqref{eq:scalar-mala-kula} daje
$\operatorname{Area}(B(p,r))
=\pi r^2(1-K(p)r^2/12+O(r^3))$, zgodnie z
obwodem~\eqref{eq:gauss-maly-okrag} po scałkowaniu.
Sam wzór na obwód otrzymujemy bez różniczkowania symbolu $O$:
dla $n=2$ całka gęstości $j(r,v)$ po $S_pM$ jest długością
okręgu geodezyjnego, więc
$C(r)=2\pi r(1-K(p)r^2/6+O(r^3))$.

Ślady tracą część danych: dodatnie i ujemne krzywizny
sekcyjne mogą się w nich znosić. W szczególności zerowy
$\mathrm{Ric}$ nie wymusza $R=0$ w wymiarze cztery.
W metrykach lorentzowskich przy śladzie występują też
znaki $\varepsilon_i=g(e_i,e_i)=\pm1$; powyższa
interpretacja objętości małych \emph{kul metrycznych}
dotyczy metryki riemannowskiej.

\begin{twierdzenie}[przestrzenie o stałej krzywiźnie]
\label{thm:stala-krzywizna}
Jeżeli wszystkie płaszczyzny styczne w rozmaitości
riemannowskiej wymiaru $n\geq2$ mają tę samą stałą
krzywiznę $\kappa$, to
\begin{equation}\label{eq:stala-krzywizna-tensor}
 R(X,Y)Z=\kappa\bigl(g(Y,Z)X-g(X,Z)Y\bigr),
 \quad \mathrm{Ric}=(n-1)\kappa g,
 \quad \mathrm{Scal}=n(n-1)\kappa.
\end{equation}
\end{twierdzenie}

\begin{proof}
Prawa strona wzoru na $R$ ma symetrie tensora
Riemanna i daje $K(\sigma)=\kappa$ dla każdej
płaszczyzny. Wykażmy, że te wartości wyznaczają
cały tensor. Niech $A$ będzie różnicą obu tensorów
z opuszczonym ostatnim indeksem. Ma ona symetrie
z Twierdzenia~\ref{thm:symetrie-riemanna} i
$A(u,v,v,u)=0$ dla wszystkich $u,v$.
Po podstawieniu $u+w$ w miejsce $u$ i użyciu
antysymetrii oraz symetrii zamiany par dostajemy
$A(u,v,v,w)=0$. Następnie podstawiamy $v+x$
w miejsce $v$: po odrzuceniu dwóch zerowych
wyrazów zostaje
\[
 A(u,v,x,w)+A(u,x,v,w)=0.
\]
Z antysymetrii w pierwszej parze wynika
$A(u,v,x,w)=A(x,u,v,w)$, a po cyklicznej
zamianie argumentów także
$A(u,v,x,w)=A(v,x,u,w)$. Pierwsza tożsamość
Bianchiego mówi, że suma tych trzech równych
wartości wynosi zero. Zatem $A=0$.
Ślad $W\mapsto R(W,U)V$ wzoru na $R$ jest
$\kappa(ng(U,V)-g(U,V))$, co daje $\mathrm{Ric}$.
Jeszcze jeden ślad daje $\mathrm{Scal}$.
\end{proof}

\begin{przyklad}[trzy geometrie modelowe]\label{prz:krzywizna-modele}
W $\R^n$ standardowe symbole Christoffela znikają:
$R=\mathrm{Ric}=\mathrm{Scal}=0$. Na sferze $S_a^n$
o promieniu $a$ w $\R^{n+1}$ operator kształtu
jest $S=\pm a^{-1}\mathrm{Id}$. Równanie Gaussa
\eqref{eq:gauss-rownanie}, ważne także dla
hiperpowierzchni, daje
$R(X,Y)Z=a^{-2}(g(Y,Z)X-g(X,Z)Y)$.
Zatem $K=a^{-2}$, $\mathrm{Ric}=(n-1)a^{-2}g$
i $\mathrm{Scal}=n(n-1)a^{-2}$.

Model hiperboloidy $\mathbb H^n=\{p:\langle p,p\rangle_L=-1,
p_0>0\}$ leży w płaskiej przestrzeni Minkowskiego.
Jej normalna $p$ ma kwadrat długości $-1$.
Z $Y\langle Z,p\rangle_L=0$ wynika
$\langle D_YZ,p\rangle_L=-g(Y,Z)$, a więc
$D_YZ=\nabla_YZ+g(Y,Z)p$. Ponadto $D_Xp=X$.
Biorąc część styczną komutatora płaskiej pochodnej
$D$ tak jak w dowodzie~\ref{thm:egregium}, uzyskujemy
\[
 R(X,Y)Z=-\bigl(g(Y,Z)X-g(X,Z)Y\bigr).
\]
Stąd $K=-1$, $\mathrm{Ric}=-(n-1)g$ i
$\mathrm{Scal}=-n(n-1)$. Po zmianie skali metryki
na $g_{\mathrm{nowe}}=a^2g$, $a>0$, otrzymujemy
$K=-a^{-2}$, $\mathrm{Ric}_{\mathrm{nowe}}
=-(n-1)a^{-2}g_{\mathrm{nowe}}=\mathrm{Ric}$ i
$\mathrm{Scal}_{\mathrm{nowe}}=-n(n-1)a^{-2}$.
Izometrie modeli z Przykładu~\ref{prz:trzy-modele-hiperboliczne}
przenoszą te wartości na model półprzestrzeni
i model kuli.
\end{przyklad}

\begin{stwierdzenie}[krzywizna w każdym z trzech modeli]
\label{prop:krzywizna-trzy-modele-hiperboliczne}
Metryki hiperboloidy, kuli Poincarégo
\eqref{eq:hip-kula} i półprzestrzeni
\eqref{eq:hip-polprzestrzen} mają tę samą
krzywiznę sekcyjną $K=-1$ w każdej płaszczyźnie.
W wymiarze $n$ w każdym modelu
$\mathrm{Ric}=-(n-1)g$ i
$\mathrm{Scal}=-n(n-1)$. Po pomnożeniu metryki
przez $a^2$ wartości te przyjmują odpowiednio
postać $K=-a^{-2}$,
$\mathrm{Ric}=-(n-1)a^{-2}g_{\mathrm{nowe}}$
oraz $\mathrm{Scal}=-n(n-1)a^{-2}$.
\end{stwierdzenie}

\begin{proof}
Dla hiperboloidy wynik otrzymaliśmy z płaskiej
pochodnej otaczającej w Przykładzie
\ref{prz:krzywizna-modele}. Odwzorowania
$P$ i $F$ z Przykładu
\ref{prz:trzy-modele-hiperboliczne} spełniają
$P^*g_{\mathcal H}=g_B$ oraz
$F^*g_{\mathbb H}=g_B$. Są zatem izometriami.
Jednoznaczność koneksji Leviego-Civity daje dla
izometrii $I$ równość
$\dd I(\nabla_XY)=\nabla_{\dd I X}(\dd I Y)$
(po odpowiednim przeniesieniu pól); wstawienie
jej do~\eqref{eq:riemann-def} przenosi tensor
$R$, a więc i~\eqref{eq:krzywizna-sekcyjna}.
To dowodzi tezy we wszystkich wymiarach.

Współrzędne $z$ pozwalają również sprawdzić
krzywiznę hiperboloidy \emph{bezpośrednio z symboli}.
Pisząc $t^2=1+|z|^2$, z
\eqref{eq:gamma-hiperboloida} mamy
\[
 \partial_i\Gamma^k_{j\ell}
   =-\delta_i^k g_{j\ell}-z^k\partial_i g_{j\ell},
 \qquad
 \Gamma^k_{ic}\Gamma^c_{j\ell}
   =\frac{z^k z_i}{t^2}g_{j\ell},
\]
ponieważ $z^c g_{ic}=z_i/t^2$. Z obliczonej
w rozdziale o koneksji pochodnej $g_{j\ell}$ wynika
\[
 \partial_i g_{j\ell}-\partial_j g_{i\ell}
    =\frac{\delta_{j\ell}z_i
                   -\delta_{i\ell}z_j}{t^2}.
\]
Po podstawieniu do~\eqref{eq:riemann-wspolrzedne}
wyrazy zawierające $z^k$ znoszą się. Zostaje
\begin{equation}\label{eq:riemann-hiperboloida-wspolrzedne}
 R^k{}_{\ell i j}
   =-\bigl(\delta_i^k g_{j\ell}
                        -\delta_j^k g_{i\ell}\bigr),
\end{equation}
co ponownie daje $K=-1$ dla każdej
płaszczyzny stycznej.

Policzmy dodatkowo krzywiznę \emph{wprost}
w wymiarze dwa. Dla metryki konforemnej
$g=e^{2\varphi}(\dd x^2+\dd y^2)$ wzór
\eqref{eq:christoffel-metryka} daje
\[
 \Gamma^k_{ij}
   =\delta_i^k\varphi_j+\delta_j^k\varphi_i
       -\delta_{ij}\varphi_k,
 \qquad \varphi_j=\partial_j\varphi.
\]
Przykładowo
$\Gamma^1_{11}=\varphi_x$,
$\Gamma^1_{12}=\varphi_y$,
$\Gamma^1_{22}=-\varphi_x$,
$\Gamma^2_{12}=\varphi_x$ i
$\Gamma^2_{22}=\varphi_y$.
W składniku $R^1{}_{2 1 2}$ wzoru
\eqref{eq:riemann-wspolrzedne} pochodne dają
$-\varphi_{xx}-\varphi_{yy}$, zaś dwa wyrazy
kwadratowe w symbolach Christoffela są
$(-\varphi_x^2+\varphi_y^2)$ i
$-(\varphi_y^2-\varphi_x^2)$, więc się znoszą.
Mianownik w~\eqref{eq:krzywizna-sekcyjna}
wynosi $e^{4\varphi}$, a licznik
$e^{2\varphi}R^1{}_{2 1 2}$; stąd
\begin{equation}\label{eq:gauss-konforemna}
 K=-e^{-2\varphi}\bigl(\varphi_{xx}
                                  +\varphi_{yy}\bigr).
\end{equation}
Dla półpłaszczyzny $\varphi=-\log y$:
$\Delta_{\mathrm E}\varphi=1/y^2$
i $e^{-2\varphi}=y^2$, co daje $K=-1$.
Dla dysku $r^2=x^2+y^2$ i
$\varphi=\log2-\log(1-r^2)$; obliczenie pochodnych
\[
 \varphi_x=\frac{2x}{1-r^2},\qquad
 \varphi_{xx}=\frac{2}{1-r^2}
                  +\frac{4x^2}{(1-r^2)^2}
\]
oraz analogiczne dla $y$ daje
$\Delta_{\mathrm E}\varphi=4/(1-r^2)^2$.
Ponieważ $e^{-2\varphi}=(1-r^2)^2/4$,
ponownie $K=-1$. Wzory na Ricciego i skalar
wynikają ze śladów w
\eqref{eq:stala-krzywizna-tensor}.

Przy skalowaniu $g\mapsto a^2g$ symbole
Christoffela pozostają takie same, więc operator
$R(X,Y)Z$ się nie zmienia. Licznik i mianownik
w~\eqref{eq:krzywizna-sekcyjna} mnożą się
odpowiednio przez $a^2$ i $a^4$; dlatego
$K\mapsto a^{-2}K$. Tensor Ricciego typu $(0,2)$
pozostaje niezmieniony: jego definicja bierze ślad
niezmienionego operatora $W\mapsto R(W,U)V$.
Dopiero ślad z nową odwrotnością metryki daje
$\mathrm{Scal}_{\mathrm{nowe}}=a^{-2}\mathrm{Scal}$.
Nie należy mylić niezmienionego tensora Ricciego
z jego wartościami na wektorach jednostkowych dla nowej metryki.
\end{proof}

\begin{przyklad}[różne płaszczyzny w jednym punkcie]
Na produkcie $S_a^2\times\R$ z metryką produktową
płaszczyzna styczna do sfery ma krzywiznę $a^{-2}$,
a płaszczyzna rozpięta przez kierunek sferyczny
i kierunek prostej ma krzywiznę $0$. W bazie
ortonormalnej złożonej z dwóch kierunków sfery
i jednego kierunku prostej tensor Ricciego ma
wartości diagonalne $(a^{-2},a^{-2},0)$,
zaś $\mathrm{Scal}=2a^{-2}$. Jedna liczba skalarna
nie mówi, która z płaszczyzn jest zakrzywiona.
\end{przyklad}

\begin{przyklad}[pomiar krzywizny obwodem koła]
\label{prz:obwod-kola-krzywizna}
Niech $\rho$ będzie \emph{odległością geodezyjną}
od ustalonego punktu na dwuwymiarowej przestrzeni
o stałej krzywiźnie. Współrzędne biegunowe na
sferze promienia $a$ dają
$g=\dd\rho^2+a^2\sin^2(\rho/a)\dd\phi^2$.
Na hiperboloidzie można zapisać punkt jako
$\bigl(a\cosh(\rho/a),
 a\sinh(\rho/a)\cos\phi,
 a\sinh(\rho/a)\sin\phi\bigr)$;
różniczkując tę parametryzację w metryce
Minkowskiego, dostajemy
$g=\dd\rho^2+a^2\sinh^2(\rho/a)\dd\phi^2$.
Na płaszczyźnie jest $g=\dd\rho^2+\rho^2\dd\phi^2$.
Okrąg $\rho=\mathrm{const}$ ma zatem obwód
$\int_0^{2\pi}\sqrt{g_{\phi\phi}}\,\dd\phi$, czyli
\begin{equation}\label{eq:obwod-kola-krzywizna}
 \begin{array}{c|c|c}
  \text{powierzchnia}& K & \text{obwód }C(\rho)\\ \hline
  S_a^2&+a^{-2}&2\pi a\sin(\rho/a),\quad 0<\rho<\pi a\\
  \R^2&0&2\pi\rho\\
  \mathbb H_a^2&-a^{-2}&2\pi a\sinh(\rho/a)
 \end{array}
\end{equation}
Ponieważ $\sin x=x-x^3/6+O(x^5)$ oraz
$\sinh x=x+x^3/6+O(x^5)$, mały okrąg na
sferze ma obwód \emph{mniejszy} od $2\pi\rho$,
a w przestrzeni hiperbolicznej \emph{większy}.
W każdym przypadku $C(\rho)=2\pi\rho
\bigl(1-K\rho^2/6+O(\rho^4)\bigr)$.
To sposób odczytania znaku krzywizny z pomiarów
wykonanych na samej powierzchni.
\end{przyklad}

\section{Schwarzschild i Kerr: próżnia nie oznacza płaskości}

W tej sekcji przyjmujemy jednostki $G=c=1$ i tę samą
konwencję znaku co w~\eqref{eq:riemann-def}. Metryki
\eqref{eq:schwarzschild} i~\eqref{eq:kerr} są
lorentzowskie. Ich $\mathrm{Ric}$ oraz $\mathrm{Scal}$
dotyczą \emph{pełnej czterowymiarowej czasoprzestrzeni};
metryki indukowane na przekrojach $t=\mathrm{const}$
mają własną, inną krzywiznę.

\begin{przyklad}[krzywizna metryki Schwarzschilda]
\label{prz:krzywizna-schwarzschild}
Dla $M>0$, w obszarze $r>2M$ i $0<\theta<\pi$, zapiszmy metrykę z
\eqref{eq:schwarzschild} jako
\[
 g=-f(r)\dd t^2+f(r)^{-1}\dd r^2
             +r^2\dd\theta^2+r^2\sin^2\theta\,\dd\phi^2,
 \qquad f(r)=1-\frac{2M}{r}.
\]
Aby zobaczyć, skąd biorą się wartości krzywizny,
podajmy obliczenie najpierw dla dowolnej gładkiej funkcji $f>0$.
Niezerowe symbole Christoffela, poza zamianą
dolnych indeksów, są następujące:
\begin{align*}
 \Gamma^t_{tr}&=\frac{f'}{2f},&
 \Gamma^r_{tt}&=\frac{ff'}2,&
 \Gamma^r_{rr}&=-\frac{f'}{2f},\\
 \Gamma^r_{\theta\theta}&=-rf,&
 \Gamma^r_{\phi\phi}&=-rf\sin^2\theta,&
 \Gamma^\theta_{r\theta}&=\Gamma^\phi_{r\phi}=\frac1r,\\
 \Gamma^\theta_{\phi\phi}&=-\sin\theta\cos\theta,&
 \Gamma^\phi_{\theta\phi}&=\cot\theta.
\end{align*}
Wstawienie ich do~\eqref{eq:riemann-wspolrzedne},
a następnie wykonanie śladu
$\mathrm{Ric}_{ij}=R^a{}_{j a i}$ daje
\begin{align}\label{eq:ricci-f-sferyczna}
 \mathrm{Ric}_{tt}&=f\left(\frac{f''}{2}+\frac{f'}r\right),
 &\mathrm{Ric}_{rr}&=-\frac1f
               \left(\frac{f''}{2}+\frac{f'}r\right),\notag\\
 \mathrm{Ric}_{\theta\theta}&=1-f-rf',
 &\mathrm{Ric}_{\phi\phi}&=(1-f-rf')\sin^2\theta,
\end{align}
zaś pozostałe składowe są zerowe. Dla kontroli
prześledźmy składową kątową. W ogólnym wzorze
\[
 \mathrm{Ric}_{ij}=\partial_a\Gamma^a_{ij}
   -\partial_j\Gamma^a_{ia}
   +\Gamma^a_{ac}\Gamma^c_{ij}
   -\Gamma^a_{jc}\Gamma^c_{ia}
\]
dla $i=j=\theta$ kolejne cztery grupy składników
wynoszą odpowiednio
\[
 -f-rf',\qquad \csc^2\theta,\qquad -2f,
 \qquad 2f-\cot^2\theta.
\]
Ich suma to $1-f-rf'$.
Teraz $f'=2M/r^2$ i $f''=-4M/r^3$:
zarówno $f''/2+f'/r$, jak i $1-f-rf'$ znikają.
Otrzymujemy
\begin{equation}\label{eq:schwarz-ricci-kretsch}
 \mathrm{Ric}=0,\qquad \mathrm{Scal}=0,
 \qquad \mathcal K:={\mathcal R}_{abcd}
                   {\mathcal R}^{abcd}=\frac{48M^2}{r^6}.
\end{equation}
Sprawdźmy również ostatnią równość, a przy okazji
wskażmy jawną składową niezerowego $R$.
Z~\eqref{eq:riemann-wspolrzedne} i podanych symboli
Christoffela znajdujemy między innymi
\[
 R^t{}_{rtr}=-\frac{f''}{2f},\qquad
 R^t{}_{\theta t\theta}=-\frac{rf'}2,\qquad
 R^r{}_{\theta r\theta}=-\frac{rf'}2,
\]
oraz, rozpisując składową kątową,
\[
 R^\theta{}_{\phi\theta\phi}
 =\partial_\theta\Gamma^\theta_{\phi\phi}
   +\Gamma^\theta_{\theta r}\Gamma^r_{\phi\phi}
   -\Gamma^\theta_{\phi\phi}
          \Gamma^\phi_{\theta\phi}
 =(1-f)\sin^2\theta.
\]
Składowe z $\phi$ w miejsce $\theta$ w parach
czasowo-kątowych i radialno-kątowych otrzymujemy
przez mnożenie przez $\sin^2\theta$. Weźmy
pseudoortonormalną ramkę o sygnaturze $(-,+,+,+)$
\[
 e_0=f^{-1/2}\partial_t,\quad
 e_1=f^{1/2}\partial_r,\quad
 e_2=r^{-1}\partial_\theta,\quad
 e_3=(r\sin\theta)^{-1}\partial_\phi
\]
i połóżmy $\mu=M/r^3$. Po obniżeniu indeksów
metryką i uwzględnieniu czynników w ramie
jedynymi niezależnymi niezerowymi wartościami są
\[
 \begin{array}{c|rrrrrr}
 (i,j)&(0,1)&(0,2)&(0,3)&(1,2)&(1,3)&(2,3)\\ \hline
 \mathcal R(e_i,e_j,e_j,e_i)
       &-2\mu&\mu&\mu&-\mu&-\mu&2\mu
 \end{array}
\]
Na przykład pierwsza wartość to $f''/2=-2\mu$,
druga to $f'/(2r)=\mu$, a ostatnia to
$(1-f)/r^2=2\mu$. Symetrie z
Twierdzenia~\ref{thm:symetrie-riemanna} dają
wszystkie pozostałe składowe; każda z sześciu
wartości występuje cztery razy przy pełnej
kontrakcji. Ponieważ w tych wyrazach indeks
czasowy występuje dwukrotnie,
\[
 \mathcal K=4\bigl((2\mu)^2+\mu^2+\mu^2
                     +\mu^2+\mu^2+(2\mu)^2\bigr)
           =48\mu^2=\frac{48M^2}{r^6}.
\]
Płaszczyzna $\operatorname{span}(\partial_\theta,
\partial_\phi)$ jest przestrzenna. W \emph{pełnej
czasoprzestrzeni} jej krzywizna sekcyjna wynosi
\begin{equation}\label{eq:schwarz-k-sekcyjna-katowa}
 K(\partial_\theta\wedge\partial_\phi)
 =\frac{r^2(1-f)\sin^2\theta}{r^4\sin^2\theta}
 =\frac{2M}{r^3}.
\end{equation}
Natomiast dwuwymiarowa sfera $t,r=\mathrm{const}$
z \emph{metryką indukowaną} ma krzywiznę Gaussa
$1/r^2$. Są to różne obiekty geometryczne;
nie wolno utożsamiać ich krzywizn. Niezmiennik
$\mathcal K$ z~\eqref{eq:schwarz-ricci-kretsch}
jest skończony przy $r=2M$, ale rośnie bez
ograniczenia, gdy $r\to0$ w przedłużeniu
czasoprzestrzeni. Zerowanie się Ricciego wynika
tu ze znoszenia się składowych śladu, a nie z
zaniku tensora Riemanna.
\end{przyklad}

\begin{przyklad}[krzywizna metryki Kerra]
\label{prz:krzywizna-kerr}
Niech $M>0$, $\Sigma=r^2+a^2\cos^2\theta$ i
$\Delta=r^2-2Mr+a^2$, jak we wzorze~\eqref{eq:kerr}.
Rachunek wykonamy dla $\Delta>0$, $\Sigma>0$ i
$0<\theta<\pi$. Wprowadźmy pseudoortonormalną ramkę
\[
 \begin{aligned}
 e_0&=\frac{(r^2+a^2)\partial_t+a\partial_\phi}
                  {\sqrt{\Sigma\Delta}},&
 e_1&=\sqrt{\frac\Delta\Sigma}\,\partial_r,\\
 e_2&=\frac{\partial_\theta}{\sqrt\Sigma},&
 e_3&=\frac{a\sin\theta\,\partial_t+
                      (\sin\theta)^{-1}\partial_\phi}{\sqrt\Sigma}.
 \end{aligned}
\]
Podstawienie do metryki daje $g(e_i,e_j)=
\operatorname{diag}(-1,1,1,1)_{ij}$. Poniższa tabela
pozwala sprawdzić oba ślady i pełną kontrakcję bez
pomijania składników mieszanych. Oznaczmy
\[
 A=\frac{Mr(r^2-3a^2\cos^2\theta)}{\Sigma^3},
 \qquad B=\frac{Ma\cos\theta(3r^2-a^2\cos^2\theta)}{\Sigma^3}.
\]
W kolejności par $01,02,03,12,13,23$ macierz
$C_{ij,kl}=\mathcal R(e_i,e_j,e_k,e_l)$ wynosi
\begin{equation}\label{eq:kerr-tabela-krzywizny}
 C=\begin{pmatrix}
 2A&0&0&0&0&-2B\\
 0&-A&0&0&-B&0\\
 0&0&-A&B&0&0\\
 0&0&B&A&0&0\\
 0&-B&0&0&A&0\\
 -2B&0&0&0&0&-2A
 \end{pmatrix}.
\end{equation}
Wszystkie inne składowe wyznaczają antysymetrie par.
Aby odtworzyć tabelę ze wzoru współrzędnego, wygodnie
położyć $x=\cos\theta$, $q=1-x^2$ i $\Sigma=r^2+a^2x^2$.
W kolejności $(t,r,x,\phi)$ niezerowe składowe odwrotnej
metryki są
\[
 \begin{aligned}
 g^{tt}&=-\frac{(r^2+a^2)^2-a^2\Delta q}{\Sigma\Delta},
 &g^{t\phi}&=-\frac{2Mar}{\Sigma\Delta},\\
 g^{\phi\phi}&=\frac{\Delta-a^2q}{\Sigma\Delta q},
 &g^{rr}&=\frac\Delta\Sigma,\qquad g^{xx}=\frac q\Sigma.
 \end{aligned}
\]
Współczynniki są wymierne w $r,x$; pochodne po $t,\phi$
znikają. Wstawiamy je do~\eqref{eq:christoffel-metryka}
i~\eqref{eq:riemann-wspolrzedne}, opuszczamy indeks metryką
i przechodzimy do powyższej ramki. Na przykład dla
$F_0=(r^2+a^2)\partial_t+a\partial_\phi$,
$F_1=\partial_r$, $F_2=-\partial_x$,
$F_3=aq\partial_t+\partial_\phi$ otrzymujemy
\[
 \begin{split}
 \mathcal R(F_0,F_1,F_0,F_1)&=
                  \frac{2Mr(r^2-3a^2x^2)}\Sigma,\\
 \mathcal R(F_0,F_1,F_2,F_3)&=
                 -\frac{2Max(3r^2-a^2x^2)}\Sigma.
 \end{split}
\]
Iloczyny czynników normalizujących w obu tych
składowych wynoszą $\Sigma^{-2}$, co daje $2A$ i $-2B$.
Pozostałe pary oblicza się tą samą skończoną procedurą.

Teraz $\mathrm{Ric}_{00}=-2A+A+A=0$,
$\mathrm{Ric}_{11}=2A-A-A=0$, a
$\mathrm{Ric}_{22}=\mathrm{Ric}_{33}=-A-A+2A=0$.
Składowe $B$ mają cztery różne indeksy, więc nie wnoszą
nic do śladu; pozadiagonalne składowe Ricciego też znikają.
Stąd $\mathrm{Ric}=0$ i $\mathrm{Scal}=0$.
Przy pełnej kontrakcji każda pozycja diagonalna tabeli
występuje cztery razy, a każda niezależna pozycja
pozadiagonalna osiem razy. Wyrazy $B$ mają jeden indeks
czasowy i wnoszą znak ujemny. Zatem
\[
 \mathcal K=4(4+1+1+1+1+4)A^2-8(4+1+1)B^2
           =48(A^2-B^2).
\]
Rozwinięcie
$r^2(r^2-3y^2)^2-y^2(3r^2-y^2)^2
=r^6-15r^4y^2+15r^2y^4-y^6$, dla $y=a\cos\theta$,
daje
\begin{equation}\label{eq:kerr-kretschmann}
 \mathcal K
 =\frac{48M^2}{\Sigma^6}
 \bigl(r^6-15a^2r^4\cos^2\theta
       +15a^4r^2\cos^4\theta-a^6\cos^6\theta\bigr).
\end{equation}
Można sprawdzić dwa ważne przypadki bez ponawiania
całego rachunku: dla $a=0$ odzyskujemy
$48M^2/r^6$ z~\eqref{eq:schwarz-ricci-kretsch},
a na równiku $\theta=\pi/2$ mamy tę samą wartość
nawet przy $a\ne0$. Poza równikiem pojawia się
zależność kątowa związana z obrotem. W sygnaturze
lorentzowskiej pełna kontrakcja $\mathcal K$ nie jest
kwadratem dodatniej normy: może mieć różne znaki.
Współrzędny rachunek krzywizny jest wymierny w $r,x$;
otrzymane tożsamości po sprowadzeniu do wspólnego mianownika
są tożsamościami wielomianowymi. Obowiązują więc również
na pozostałych obszarach $\Delta\ne0$, $\Sigma>0$,
choć użyta ramka z pierwiastkiem $\Delta$ nie jest tam rzeczywista.
Równanie $\Sigma=0$ dla $a\ne0$ odpowiada
$r=0,\theta=\pi/2$ i wskazuje osobliwość
pierścieniową: wzdłuż równika $\mathcal K=48M^2/r^6$
rośnie bez ograniczenia. Usuwalność osobliwości współrzędnych
przy $\Delta=0$, $\Sigma>0$, uzasadnia inny rachunek,
nie sama skończoność $\mathcal K$. Wprowadzamy lokalnie
\[
 \dd v=\dd t+\frac{r^2+a^2}{\Delta}\dd r,
 \qquad \dd\psi=\dd\phi+\frac a\Delta\dd r.
\]
Po podstawieniu do metryki wszystkie mianowniki $\Delta$
znoszą się i otrzymujemy
\[
 \begin{split}
 g={}&-\left(1-\frac{2Mr}{\Sigma}\right)\dd v^2
       +2\dd v\dd r-\frac{4Mar\sin^2\theta}{\Sigma}\dd v\dd\psi
       -2a\sin^2\theta\,\dd r\dd\psi\\
    &+\Sigma\dd\theta^2
       +\sin^2\theta\left(r^2+a^2+
                 \frac{2Mra^2\sin^2\theta}{\Sigma}\right)\dd\psi^2.
 \end{split}
\]
Jej wyznacznik wynosi $-\Sigma^2\sin^2\theta$.
Zatem dla $0<\theta<\pi$ i $\Sigma>0$ nowy wzór jest
gładki i niezdegenerowany także przy $\Delta=0$.
Definiuje lokalne przedłużenie przez tę powierzchnię.
Te stwierdzenia dotyczą pełnej
metryki lorentzowskiej. Przekrój $t=\mathrm{const}$
może mieć niezerowy tensor Ricciego własnej
metryki trójwymiarowej.
\end{przyklad}

\begin{uwaga}[interpretacja porównawcza]
Na sferze już Ricci i skalar rejestrują dodatnią
krzywiznę. W próżniowych przykładach Schwarzschilda
i Kerra oba te ślady znikają, choć transport po
pętlach i zachowanie sąsiednich geodezyjnych ujawniają
niezerowy tensor $R$. W czterowymiarowej geometrii
lorentzowskiej pozostaje wtedy część krzywizny,
która nie jest opisana przez Ricciego (tensor Weyla).
Pełny rozkład na część bezśladową i składniki
wyznaczone przez Ricciego wykracza poza ten rozdział.
\end{uwaga}

\section{Kiedy metryka jest naprawdę płaska?}

\begin{twierdzenie}[lokalna charakteryzacja płaskości]
\label{thm:lokalna-plaskosc}
Niech $(M,g)$ będzie rozmaitością riemannowską
i $p\in M$. Następujące warunki są równoważne
w odpowiednio małym otoczeniu $p$:
\begin{enumerate}[label=\textup{(\roman*)}]
 \item tensor Riemanna znika: $R\equiv0$;
 \item transport równoległy jest lokalnie niezależny
       od drogi o ustalonych końcach;
 \item istnieje ortonormalna rama pól równoległych
       $E_1,\ldots,E_n$;
 \item istnieją współrzędne $(x^1,\ldots,x^n)$,
       w których $g=\sum_i(\dd x^i)^2$;
 \item otoczenie $p$ jest izometryczne z otwartym
       podzbiorem $\R^n$.
\end{enumerate}
Ponadto $R=0$ wtedy i tylko wtedy, gdy
$K_p(\sigma)=0$ dla każdej płaszczyzny stycznej
we wszystkich punktach tego otoczenia.
W przypadku pseudoriemannowskim równoważność
warunków \textup{(i)--(iii)} i wersji \textup{(iv)}
z metryką stałą właściwej sygnatury jest taka sama.
\end{twierdzenie}

\begin{proof}
Wybierzmy małą, ściągalną kulę współrzędnych $U$
wokół $p$. Załóżmy (i). Pokażemy dokładnie,
dlaczego transport po różnych drogach daje to samo.
Niech $H(s,t)$ będzie gładką homotopią dróg
od $p$ do ustalonego $q\in U$, z nieruchomymi
końcami, $0\leq s,t\leq1$. Przenieśmy
$v\in T_pM$ równolegle wzdłuż każdej drogi
$t\mapsto H(s,t)$; oznaczmy otrzymane pole
wzdłuż $H$ przez $V(s,t)$. Mamy $\nabla_tV=0$.
Także gdy $H$ nie jest immersją, pochodne pól wzdłuż
$H$ spełniają regułę komutacji. Współrzędnie mają
postać $D_\alpha=\partial_\alpha+A_\alpha$, gdzie
$A_\alpha=\Gamma_i(H)\partial_\alpha H^i$.
Rozwinięcie komutatora daje
$[D_t,D_s]=\partial_tA_s-\partial_sA_t+[A_t,A_s]
=R(\partial_tH,\partial_sH)$: mieszane drugie pochodne
$H$ znoszą się, a pozostałe składniki są dokładnie
wzorem~\eqref{eq:riemann-wspolrzedne}. Stąd
\[
 \nabla_t(\nabla_sV)
 =\nabla_s(\nabla_tV)+R(\partial_tH,\partial_sH)V=0.
\]
Na początku $V(s,0)=v$, stąd $\nabla_sV(s,0)=0$.
Jednoznaczność równania transportu równoległego
wzdłuż $t$ daje $\nabla_sV(s,t)=0$ dla każdego
$t$, zwłaszcza przy $t=1$. Transport jest więc
niezależny od homotopii z ustalonymi końcami.
W wypukłym obrazie kuli współrzędnych dwie drogi
łączymy homotopią liniową
$H(s,t)=x^{-1}((1-s)x(\gamma_0(t))+s x(\gamma_1(t)))$.
Zachowuje ona końce. Dla dróg odcinkami gładkich
wybieramy wspólny skończony podział parametru $t$;
powyższy rachunek stosujemy kolejno na jego przedziałach.
Na łączeniach wartości $V$ i $\nabla_sV$ są zgodne,
więc argument przechodzi przez cały podział. Dowodzi to (ii).

Wybierzmy ortonormalną bazę $e_1,\ldots,e_n$ w $T_pM$
i przenieśmy każdy $e_i$ do $q\in U$. Niezależność
od drogi definiuje pole $E_i(q)$; transport po
radialnych drogach zależy gładko od końca,
więc $E_i$ są gładkie. Idąc najpierw do $q$,
a potem po dowolnej krótkiej krzywej przez $q$,
widzimy $\nabla_XE_i=0$ dla wszystkich $X$.
Zgodność koneksji z metryką zachowuje iloczyny
skalarne: $g(E_i,E_j)=\delta_{ij}$ na $U$.
Otrzymaliśmy (iii).

Koneksja Leviego-Civity nie ma torsji, więc
\[
 [E_i,E_j]=\nabla_{E_i}E_j-\nabla_{E_j}E_i=0.
\]
Przepływy komutujących pól komutują lokalnie.
Odwzorowanie otrzymane przez wykonanie po kolei
ich przepływów z czasami $x^1,\ldots,x^n$,
zaczynając od $p$, jest lokalnym dyfeomorfizmem
(jego różniczka w zerze wysyła bazę standardową
na $E_i(p)$). Komutowanie przepływów sprawia,
że w powstałych współrzędnych $E_i=\partial_i$.
Stąd $g_{ij}=g(E_i,E_j)=\delta_{ij}$, czyli (iv).
Odwzorowanie współrzędnych jest wtedy izometrią
na otwarty podzbiór $\R^n$, więc (iv) daje (v).
Izometria zachowuje koneksję i $R$, a w $\R^n$
standardowa koneksja ma $R=0$: (v) daje (i).
Implikacje (iii)$\Rightarrow$(ii) można też
sprawdzić bezpośrednio: współczynniki wektora
w ramie równoległej nie zmieniają się podczas
transportu.

Jeśli $R=0$, wzór~\eqref{eq:krzywizna-sekcyjna}
daje $K=0$. Odwrotność jest dokładnie krokiem
polaryzacji w dowodzie Twierdzenia
\ref{thm:stala-krzywizna} dla $\kappa=0$.
W sygnaturze pseudoriemannowskiej wybieramy
bazę ortonormalną z $g(e_i,e_i)=\pm1$; reszta
dowodu pozostaje bez zmian, a metryka w
współrzędnych ma stałe współczynniki $\pm1$.
\end{proof}

\begin{wniosek}[niski wymiar a płaskość]
\label{cor:ricci-plaskosc-3d}
Na powierzchni riemannowskiej $\mathrm{Scal}=0$
w otoczeniu wtedy i tylko wtedy, gdy metryka jest
tam lokalnie płaska. W wymiarze trzy wystarcza
$\mathrm{Ric}=0$ w otoczeniu. W wymiarze cztery
dla metryk lorentzowskich $\mathrm{Ric}=0$
nie wymusza już płaskości.
\end{wniosek}
\begin{proof}
W wymiarze dwa $\mathrm{Scal}=2K$ na mocy
Stwierdzenia~\ref{prop:ricci-sekcje}; stosujemy
Twierdzenie~\ref{thm:lokalna-plaskosc}.
W wymiarze trzy i przy $\mathrm{Ric}=0$ wybierzmy
\emph{dowolną} bazę ortonormalną $e_1,e_2,e_3$.
Oznaczając $K_{ij}=K(e_i\wedge e_j)$, dostajemy
z~\eqref{eq:ricci-suma-sekcji}
\[
 K_{12}+K_{13}=0,\qquad
 K_{12}+K_{23}=0,\qquad
 K_{13}+K_{23}=0.
\]
Po odjęciu dwóch pierwszych równań mamy
$K_{13}=K_{23}$; trzecie wymusza
$K_{13}=K_{23}=0$, a pierwsze daje $K_{12}=0$.
Każda płaszczyzna w $T_pM$ może być pierwszą
płaszczyzną pewnej bazy ortonormalnej; skoro
baza była dowolna, wszystkie krzywizny sekcyjne
znikają. Twierdzenie~\ref{thm:lokalna-plaskosc}
daje płaskość. Odwrotne implikacje są oczywiste,
bo $R=0$ pociąga oba zerowe ślady.
Kontrprzykłady w wymiarze cztery dają
\eqref{eq:schwarz-ricci-kretsch} i
\eqref{eq:kerr-kretschmann}: są lorentzowskie
i mają $\mathrm{Ric}=0$ przy $R\ne0$.
\end{proof}

\begin{uwaga}[lokalnie a globalnie]
Twierdzenie jest lokalne. Płaski torus
$\R^2/\mathbb Z^2$ ma $R=0$ i lokalnie wygląda
jak płaszczyzna, lecz globalnie ma inną topologię
i nie jest izometryczny z całą $\R^2$.
Walec z Przykładu~\ref{prz:egregium-cztery}
również jest lokalnie płaski. Sama równość
$\mathrm{Scal}=0$ jest znacznie słabsza od
$R=0$, co pokazują przykłady Schwarzschilda
i Kerra.
\end{uwaga}

\section{Druga wariacja energii i pola Jacobiego}
\label{sec:druga-wariacja-jacobi}

Pierwsza wariacja z
Twierdzenia~\ref{thm:pierwsza-wariacja-dlugosci}
mówi, \emph{gdzie} długość może mieć minimum:
krzywa minimalna, po parametryzacji długością łuku,
jest geodezyjną. Drugie różniczkowanie mówi,
\emph{czy w pobliżu geodezyjnej istnieje kierunek},
w którym energia maleje. Krzywizna pojawia się
tu dlatego, że zamiana kolejności dwóch pochodnych
\emph{kowariantnych} wzdłuż wariacji nie jest
na ogół dozwolona bez poprawki $R$.

W całej sekcji $(M,g)$ jest rozmaitością
\emph{riemannowską}, $\nabla$ jej koneksją
Leviego-Civity, a znak $R$ jest ustalony
w~\eqref{eq:riemann-def}. Nie przenosimy
wniosków o minimizowaniu odległości na całe
czasoprzestrzenie lorentzowskie.

\subsection*{Drugie różniczkowanie energii}

Niech $\gamma:[0,\ell]\to M$ będzie geodezyjną
o prędkości $1$, gdzie $\ell>0$, więc $\ell=L_g(\gamma)$.
Rozważmy gładką rodzinę $F(s,t)$ krzywych
o stałych końcach
$F(s,0)=p$, $F(s,\ell)=q$,
$F(0,t)=\gamma(t)$, i jej pole wariacyjne
$V(t)=\partial_sF(0,t)$. W szczególności
$V(0)=V(\ell)=0$. Piszemy $D_t$ oraz $D_s$
na pochodne kowariantne pól wzdłuż powierzchni
parametrów $F$.

\begin{twierdzenie}[Druga wariacja energii]
\label{thm:druga-wariacja-energii}
Dla powyższej wariacji, przy energii
z~\eqref{eq:dlugosc-energia}, zachodzi
\begin{equation}\label{eq:druga-wariacja-energii}
 \left.\frac{\dd^2}{\dd s^2}E_g(F_s)\right|_{s=0}
 =\int_0^\ell\left(
   |D_tV|_g^2-g(R(V,T)T,V)\right)\dd t,
 \qquad T=\dot\gamma.
\end{equation}
Prawą stronę oznaczamy przez $I(V,V)$.
Jej postać dwuliniowa, zwana \emph{formą indeksową}, to
\begin{equation}\label{eq:forma-indeksowa}
 I(V,W)=\int_0^\ell\left(
   g(D_tV,D_tW)-g(R(V,T)T,W)\right)\dd t
\end{equation}
dla pól $V,W$ zerujących się na końcach.
\end{twierdzenie}
\begin{proof}
Oznaczmy na całej powierzchni wariacji
$S=\partial_sF$ i $U=\partial_tF$.
Zgodność koneksji z metryką daje, dla każdego $s$,
\[
 \frac{\dd}{\dd s}E_g(F_s)
 =\int_0^\ell g(D_sU,U)\,\dd t.
\]
Ponieważ $[\partial_s,\partial_t]=0$ i koneksja
nie ma torsji, $D_sU=D_tS$. Po całkowaniu
przez części otrzymujemy \emph{dla całej rodziny}
\begin{equation}\label{eq:pierwsza-wariacja-energii-kow}
 \frac{\dd}{\dd s}E_g(F_s)
 =\bigl[g(S,U)\bigr]_0^\ell
   -\int_0^\ell g(S,D_tU)\,\dd t.
\end{equation}
Końce są nieruchome, więc $S(s,0)=S(s,\ell)=0$;
składnik brzegowy oraz jego pochodna po $s$
znikają. Różniczkujemy pozostałą całkę.
Przy $s=0$ mamy $D_tU=D_tT=0$, gdyż
$\gamma$ jest geodezyjną. Człon zawierający
$D_sS$ znika, bez względu na \emph{przyspieszenie}
wybranej wariacji. Zostaje
\[
 E_g''(0)=-\int_0^\ell
         g(V,D_sD_tU|_{s=0})\,\dd t.
\]
Definicja~\eqref{eq:riemann-def} na polach
wzdłuż $F$ ma tutaj postać
$D_sD_tU-D_tD_sU=R(S,U)U$.
Ponieważ $D_sU=D_tS$, dostajemy krok po kroku
\[
 D_sD_tU=D_t(D_sU)+R(S,U)U
          =D_t^2S+R(S,U)U.
\]
Po podstawieniu $s=0$ jest to
$D_t^2V+R(V,T)T$. Wreszcie reguła iloczynu
i $V(0)=V(\ell)=0$ dają
\[
 -\int_0^\ell g(V,D_t^2V)\,\dd t
 =-\bigl[g(V,D_tV)\bigr]_0^\ell
  +\int_0^\ell|D_tV|_g^2\,\dd t
 =\int_0^\ell|D_tV|_g^2\,\dd t.
\]
To dowodzi~\eqref{eq:druga-wariacja-energii}.
Polaryzacja daje~\eqref{eq:forma-indeksowa};
symetria $I(V,W)=I(W,V)$ wynika również
z symetrii tensora krzywizny, bo
$g(R(V,T)T,W)=g(R(W,T)T,V)$.
\end{proof}

W przeciwieństwie do pierwszej wariacji,
w~\eqref{eq:druga-wariacja-energii} pojawia się
zarówno koszt zmieniania kierunku pola $V$,
czyli $|D_tV|^2$, jak i wkład krzywizny.
Na sferze dodatnia krzywizna odejmuje od
pierwszego składnika; przy ujemnej krzywiźnie
poprzeczne wychylenie kosztuje dodatkową energię.

\begin{figure}[htbp]
\centering
\begin{tikzpicture}[>=Latex,scale=1]
  \fill[manifoldblue!6,rounded corners=10pt]
       (-3.8,-1.45) rectangle (3.8,1.48);
  \draw[manifoldblue!45,rounded corners=10pt]
       (-3.8,-1.45) rectangle (3.8,1.48);
  \coordinate (p) at (-3.08,-.54);
  \coordinate (q) at (3.08,-.54);
  \draw[manifoldblue!50,thick]
       (p) .. controls (-1.9,1.07) and (1.85,1.04) .. (q);
  \draw[manifoldblue!50,thick]
       (p) .. controls (-1.85,-1.24) and (1.9,-1.25) .. (q);
  \draw[manifoldred,ultra thick,->]
       (p) .. controls (-1.9,.10) and (1.90,.10) .. (q);
  \draw[manifoldgreen!75!black,thick,->]
       (-1.38,-.19)--(-1.38,.28);
  \draw[manifoldgreen!75!black,thick,->]
       (0,-.03)--(0,.56);
  \draw[manifoldgreen!75!black,thick,->]
       (1.38,-.19)--(1.38,.27);
  \node[manifoldgreen!55!black,right] at (0,.49) {$V(t)$};
  \fill[manifoldblue!80!black] (p) circle (2.6pt)
       node[below left] {$p$};
  \fill[manifoldblue!80!black] (q) circle (2.6pt)
       node[below right] {$q$};
  \node[manifoldred!85!black] at (0,-.57) {$\gamma=F_0$};
  \node[manifoldblue!75!black] at (2.2,.92) {$F_s$};
\end{tikzpicture}
\caption{Wariacja o nieruchomych końcach: pole $V$ opisuje pierwsze przemieszczenie punktów krzywej $\gamma$. Druga wariacja mierzy zmianę energii przy takim wychyleniu. Rysunek schematyczny; otaczające krzywe nie muszą być geodezyjnymi.}
\label{fig:druga-wariacja-krzywej}
\end{figure}

\begin{stwierdzenie}[Druga wariacja długości]
\label{prop:druga-wariacja-dlugosci}
Przy tych samych założeniach, dla wariacji
regularnych krzywych, zachodzi
\begin{equation}\label{eq:druga-wariacja-dlugosci}
 \left.\frac{\dd^2}{\dd s^2}L_g(F_s)\right|_{s=0}
 = I(V,V)-\int_0^\ell
      g(D_tV,T)^2\,\dd t.
\end{equation}
Jeśli $V$ jest prostopadłe do $T$ w każdym
punkcie, obie drugie wariacje są równe.
\end{stwierdzenie}
\begin{proof}
Przy $s=0$ mamy $|U|=|T|=1$. Zwykłe
różniczkowanie funkcji $\sqrt{h}$ dla
$h(s,t)=g(U,U)$ pokazuje punktowo, że
\[
 \left.\partial_s^2|U|\right|_{s=0}
 =\left.\partial_s^2\frac{|U|^2}{2}\right|_{s=0}
   -\bigl(\left.\partial_s|U|\right|_{s=0}\bigr)^2.
\]
Pierwszy składnik po scałkowaniu daje $E''(0)$.
Ponadto $\partial_s|U||_0=g(D_sU,T)|_0
=g(D_tV,T)$. To daje wzór. Jeżeli
$g(V,T)=0$ na całym $[0,\ell]$, to
$g(D_tV,T)=\frac{\dd}{\dd t}g(V,T)=0$,
bo $D_tT=0$.
\end{proof}

To rozróżnienie jest potrzebne: dla pola
$V=fT$ z $f(0)=f(\ell)=0$ wzór na energię
daje $I(fT,fT)=\int_0^\ell(f')^2\dd t$,
natomiast druga wariacja długości wynosi
zero. Przesuwanie punktów \emph{wzdłuż} tej
samej krzywej zmienia jej parametryzację,
a długość nie zależy od parametryzacji.

\subsection*{Pola Jacobiego: sąsiednie geodezyjne}

\begin{definicja}[Pole Jacobiego]
\index{pole jacobiego@Pole Jacobiego}
\label{def:pole-jacobiego}
Pole wektorowe $J$ wzdłuż geodezyjnej
$\gamma$, z $T=\dot\gamma$, nazywamy
\emph{polem Jacobiego}, jeśli
\begin{equation}\label{eq:jacobi-rownanie}
 D_t^2J+R(J,T)T=0.
\end{equation}
Równanie jest liniowe. Przy naszej konwencji
na sferze $R(J,T)T$ ma zwrot prostopadłego
pola $J$, a więc $D_t^2J=-R(J,T)T$
jest skierowane z powrotem ku geodezyjnej.
\end{definicja}

\begin{twierdzenie}[Pola Jacobiego powstają z geodezyjnych]
\label{thm:jacobi-wariacja-geodezyjnych}
Jeżeli $F(s,t)$ jest wariacją, dla której
każda krzywa $t\mapsto F(s,t)$ jest
geodezyjną z parametrem afinicznym, to
$J=\partial_sF|_{s=0}$ spełnia
\eqref{eq:jacobi-rownanie}. Odwrotnie:
każde pole Jacobiego wzdłuż $\gamma$ na
zwartym przedziale jest polem wariacyjnym
pewnej rodziny geodezyjnych.
\end{twierdzenie}
\begin{proof}
Niech znów $S=\partial_sF$ i $U=\partial_tF$.
Teraz $D_tU=0$ dla \emph{każdego} $s$,
więc $D_sD_tU=0$. Obliczenie wykonane
w dowodzie Twierdzenia~\ref{thm:druga-wariacja-energii}
daje
\[
 0=D_t^2S+R(S,U)U.
\]
Przy $s=0$ jest to równanie Jacobiego.

Dla odwrotności wybierzmy
$A=J(0)$ i $B=D_tJ(0)$.
Weźmy krzywą $p_s$ z $p_0=\gamma(0)$,
$\dot p_0=A$. Niech $P_s:T_{p_0}M\to T_{p_s}M$
będzie transportem równoległym wzdłuż $p_s$
i połóżmy $v_s=P_s(T(0)+sB)$.
Dla małych $s$ rodzina geodezyjnych
$F(s,t)=\gamma_{p_s,v_s}(t)$ istnieje na
całym zadanym zwartym przedziale: rozwiązania
ODE zależą gładko od danych bliskich danym
$\gamma$, a skończona liczba lokalnych
przedłużeń pokrywa przedział.
Jej pole $\widetilde J=\partial_sF|_0$
ma $\widetilde J(0)=A$ oraz, z
$D_t\widetilde J=D_sU$,
$D_t\widetilde J(0)=D_s v_s|_0=B$.
W równoległej ramie wzdłuż $\gamma$
równanie~\eqref{eq:jacobi-rownanie} staje się
układem liniowych ODE drugiego rzędu
$y''+A(t)y=0$. Jednoznaczność rozwiązania
przy ustalonych $J(0),D_tJ(0)$ daje
$\widetilde J=J$ na całym przedziale.
\end{proof}

Pole Jacobiego mierzy więc, jak zmienia się
położenie punktu na geodezyjnej, gdy nieznacznie
zmieniamy jej \emph{punkt startowy lub prędkość}.
Nie zakładamy tu, że końce są nieruchome;
wariacje geodezyjnych o ustalonych obu końcach
mogą nie istnieć poza szczególnymi sytuacjami.

\begin{definicja}[Punkty sprzężone]
\index{punkty sprzezone@Punkty sprzężone}
\label{def:punkty-sprzezone}
Parametry $0$ i $c>0$, a także odpowiadające
im punkty $p=\gamma(0)$ i $\gamma(c)$
\emph{wzdłuż wybranej geodezyjnej}, nazywamy
\emph{sprzężonymi}, jeśli istnieje niezerowe
pole Jacobiego $J$ spełniające
$J(0)=J(c)=0$. Gdy geodezyjna sama się
przecina, sformułowanie przez parametry
usuwa niejednoznaczność.
\end{definicja}

\begin{stwierdzenie}[Punkty sprzężone i odwzorowanie wykładnicze]
\label{prop:jacobi-dexp}
Jeżeli $\gamma(t)=\exp_p(tv)$ oraz $v=T(0)$,
to dla $w\in T_pM$ pole powstające przez
zmianę prędkości początkowej $v+sw$ ma postać
\begin{equation}\label{eq:jacobi-dexp}
 J_w(t)=\left.\partial_s\exp_p(t(v+sw))\right|_{s=0}
       =\dd(\exp_p)_{tv}(tw),
 \qquad J_w(0)=0,\quad D_tJ_w(0)=w.
\end{equation}
W konsekwencji $\gamma(c)$ jest sprzężony z $p$
wzdłuż $\gamma$ wtedy i tylko wtedy, gdy
$\dd(\exp_p)_{cv}$ ma niezerowe jądro.
\end{stwierdzenie}
\begin{proof}
Różniczkujemy rodzinę geodezyjnych
$t\mapsto\exp_p(t(v+sw))$ po $s$ i stosujemy
regułę łańcuchową; daje to pierwsze dwie
równości. Wszystkie krzywe startują z $p$,
więc $J_w(0)=0$. Ich prędkości początkowe
są równe $v+sw$, stąd $D_tJ_w(0)=w$.
Każde pole Jacobiego zerujące się w $0$
jest jednym z $J_w$, ponieważ jego
początkowa pochodna $w$ wyznacza je jednoznacznie.
Dla $c>0$ warunek $J_w(c)=0$ jest równoważny
$\dd(\exp_p)_{cv}(cw)=0$. Ponieważ
$w\ne0$ dokładnie wtedy, gdy $cw\ne0$,
otrzymujemy tezę.
\end{proof}

Różniczka $\dd\exp_p$ pokazuje zatem,
kiedy geodezyjne startujące z $p$ zaczynają
się \emph{ogniskować}: początkowo różne kierunki
mogą w pierwszym przybliżeniu prowadzić
do tego samego punktu. W punkcie sprzężonym
normalne współrzędne z centrum $p$ tracą
lokalną odwracalność wzdłuż tego promienia.

\subsection*{Trzy znaki krzywizny i przykłady}

Na przestrzeni o stałej krzywiźnie sekcyjnej
$\kappa$, z~\eqref{eq:stala-krzywizna-tensor} wynika,
że dla jednostkowego $T$ i prostopadłego
równoległego pola $E$ mamy
$R(E,T)T=\kappa E$.
Jeżeli $J(t)=f(t)E(t)$, równanie Jacobiego
redukuje się do zwykłego równania
\begin{equation}\label{eq:jacobi-stala-krzywizna}
 f''+\kappa f=0,
 \qquad f(0)=0,\quad f'(0)=1.
\end{equation}
Jego rozwiązania dla trzech znaków krzywizny to
\begin{equation}\label{eq:jacobi-trzy-krzywizny}
 f_\kappa(t)=
 \begin{cases}
  \dfrac{\sin(\sqrt\kappa\,t)}{\sqrt\kappa},
       &\kappa>0,\\[5pt]
  t,&\kappa=0,\\[3pt]
  \dfrac{\sinh(\sqrt{-\kappa}\,t)}{\sqrt{-\kappa}},
       &\kappa<0.
 \end{cases}
\end{equation}
Sprawdzamy to bez skrótu: dla $\kappa>0$
dwie pochodne sinusa dają $f''=-\kappa f$;
analogicznie $\sinh$ daje $f''=(-\kappa)f$,
więc $f''+\kappa f=0$; wszystkie trzy wzory
mają $f(0)=0$, $f'(0)=1$.
Składowa styczna pola Jacobiego ma natomiast
postać $(a+bt)T$, bo $R(T,T)T=0$.
Niezerowa taka składowa nie może znikać
w dwóch różnych chwilach. O punktach
sprzężonych decydują więc składowe poprzeczne.

\begin{przyklad}[Sfera o promieniu $R$]
\label{prz:jacobi-sfera}
Na $S_R^n$, $n\geq2$ i $R>0$, mamy $\kappa=1/R^2$.
Dla $0\ne w\perp v$, po równoległym przesunięciu
$w$ wzdłuż jednostkowej geodezyjnej, dostajemy
\[
 J_w(t)=R\sin(t/R)\,P_t w.
\]
Pole startuje od zera, rośnie, a następnie
znika ponownie przy $t=\pi R$, czyli w antypodzie.
Istnieje $n-1$ niezależnych kierunków $w$,
więc jest to punkt sprzężony krotności $n-1$.
Geometrycznie: różne wielkie półokręgi
wychodzą z $p$ i spotykają się w $-p$.
Na łuku dłuższym niż $\pi R$ można zobaczyć
utratę minimalności bez żadnego ogólnego
twierdzenia. Dla $\ell>\pi R$ weźmy
równoległe jednostkowe pole normalne $E$
i $V(t)=\sin(\pi t/\ell)E(t)$.
Wtedy $V(0)=V(\ell)=0$ i, z formy indeksowej,
\begin{equation}\label{eq:jacobi-sfera-indeks}
 I(V,V)
 =\int_0^\ell\left[\frac{\pi^2}{\ell^2}
                   \cos^2\!\left(\frac{\pi t}{\ell}\right)
       -\frac1{R^2}\sin^2\!\left(\frac{\pi t}{\ell}\right)
                    \right]\dd t
 =\frac\ell2\left(\frac{\pi^2}{\ell^2}
                     -\frac1{R^2}\right)<0.
\end{equation}
Taki łuk nie może być minimalny.
Przy $\ell=\pi R$ wybrane pole daje $I=0$,
a mimo to półokrąg \emph{jest} minimalny:
zero drugiej wariacji samo nie rozstrzyga sprawy.
\end{przyklad}

\begin{figure}[htbp]
\centering
\begin{tikzpicture}[>=Latex]
 \shade[inner color=white,outer color=manifoldblue!46] (0,0) circle (1.72);
 \draw[manifoldblue!75!black,very thick] (0,0) circle (1.72);
 \draw[manifoldblue!40] (-1.72,0)--(1.72,0);
 \foreach \b in {-15,50}
  \draw[manifoldblue!65,thick] plot[domain=0:180,samples=90]
    ({1.72*sin(\b)*sin(\x)},{1.72*cos(\x)});
 \draw[manifoldred,ultra thick,postaction={decorate,decoration={markings,
   mark=at position .34 with {\arrow{Latex}}}}]
  plot[domain=0:180,samples=90] ({1.72*sin(20)*sin(\x)},{1.72*cos(\x)});
 \draw[manifoldgreen!75!black,very thick,->]
  ({1.72*sin(20)},0)--({1.72*sin(20)+.3*1.72*cos(20)},0)
  node[midway,above] {$J$};
 \fill[manifoldblue!85!black] (0,1.72) circle (2.4pt) node[above] {$p$};
 \fill[manifoldblue!85!black] (0,-1.72) circle (2.4pt) node[below] {$-p$};
 \node[manifoldred!85!black] at (.92,.95) {$\gamma$};
 \node[manifoldblue!70!black,anchor=west] at (2,-.5) {$J(0)=J(\pi R)=0$};
\end{tikzpicture}
\caption{Rzut wielkich półokręgów na sferze oglądanej prostopadle do osi biegunowej.
Półokręgi spotykają się w antypodzie. Pole Jacobiego, otrzymane przez zmianę
długości geograficznej, znika przy $p$ i $-p$. Zielona strzałka pokazuje jego
rzut w punkcie równika, w umownej skali.}
\label{fig:jacobi-sfera-ogniskowanie}
\end{figure}

\begin{przyklad}[Płaszczyzna i trzy modele hiperboliczne]
\label{prz:jacobi-plaska-hiperboliczna}
W $\R^n$ pole poprzeczne z $J(0)=0$
i $D_tJ(0)=w$ ma postać $J(t)=tw$,
przy identyfikacji przestrzeni stycznych przez
transport równoległy. Dla $w\ne0$ nie znika już przy $t>0$.
W przestrzeni o $K=-1$ analogicznie
$J(t)=\sinh(t)P_tw$, więc sąsiednie geodezyjne
rozchodzą się szybciej. Ponadto dla
$V=V^\parallel+V^\perp$ z zerowymi końcami
\[
 I(V,V)=\int_0^\ell
        \bigl(|D_tV|^2+|V^\perp|^2\bigr)\dd t>0
 \quad\text{dla }V\ne0.
\]
Nie ma punktów sprzężonych. Dotyczy to
\emph{tej samej} przestrzeni hiperbolicznej
zapisanej jako półprzestrzeń, kula Poincarégo
lub hiperboloida; metryki i izometrie między
nimi podano w~\ref{prz:trzy-modele-hiperboliczne},
a $K=-1$ obliczono w
Stwierdzeniu~\ref{prop:krzywizna-trzy-modele-hiperboliczne}.
Wzór dodatniej formy mówi o \emph{lokalnej}
stabilności. Globalną minimalność i jedyność
geodezyjnych w tych trzech modelach potwierdzają
odległości~\eqref{eq:odleglosci-modele-hiperboliczne}.
\end{przyklad}

\subsection*{Co naprawdę wynika dla krzywej minimalnej?}

\begin{twierdzenie}[Warunki konieczne i punkty sprzężone]
\label{thm:minimalna-bez-sprzezonego-wewnatrz}
Jeżeli ciągła, odcinkami $C^1$ krzywa
$\gamma:[0,\ell]\to M$, $\ell>0$, minimalizuje
długość między $p$ i $q$ i jest parametryzowana
długością łuku, to:
\begin{enumerate}[label=\textup{(\roman*)}]
 \item $\gamma$ jest geodezyjną;
 \item $I(V,V)\geq0$ dla każdego gładkiego pola
       $V$ wzdłuż $\gamma$ z $V(0)=V(\ell)=0$;
 \item żaden punkt $\gamma(c)$ dla
       $0<c<\ell$ nie jest sprzężony z $p$
       wzdłuż $\gamma$.
\end{enumerate}
Punkt końcowy $q$ \emph{może} być sprzężony z $p$.
\end{twierdzenie}
\begin{proof}
Każdy podłuk $\gamma$ również minimalizuje długość:
jego skrócenie skróciłoby całą drogę. Ustalmy punkt
$\gamma(t_0)$ i jego małe otoczenie normalne. Dla
$t$ bliskich $t_0$ krzywa pozostaje w tym otoczeniu.
Warunek równości w dowodzie lokalnej minimalności
promieni (Twierdzenie~\ref{thm:lokalna-minimalnosc-geodezyjnej})
wymusza, aby każdy z dwóch podłuków wychodzących
z $\gamma(t_0)$ był radialny. Parametr długości łuku
nadaje im prędkość $1$. Są więc gładkimi geodezyjnymi
po obu stronach $t_0$, a argument o braku narożników
z dowodu Twierdzenia~\ref{thm:pierwsza-wariacja-dlugosci}
zrównuje ich prędkości w $t_0$. Jednoznaczność ODE
skleja je w jedną gładką geodezyjną. Na końcach stosujemy
argument jednostronny. To dowodzi (i), bez zakładania
gładkości minimalnej krzywej z góry.
Sama druga wariacja nie zastępuje tego kroku:
stosujemy jej wzór dopiero \emph{na geodezyjnej}.

Niech $F_s$ będzie dowolną wariacją z nieruchomymi
końcami. Dla krzywych określonych na
$[0,\ell]$ nierówność~\eqref{eq:energia-cauchy}
oraz minimalność $\gamma$ dają
\[
 E_g(F_s)\geq\frac{L_g(F_s)^2}{2\ell}
             \geq\frac{\ell^2}{2\ell}
             =E_g(\gamma).
\]
Energia ma więc minimum przy $s=0$, skąd
$E_g''(0)\geq0$. Każde gładkie pole $V$
zerujące się na końcach można zrealizować
przez $F(s,t)=\exp_{\gamma(t)}(sV(t))$
przy dostatecznie małym $s$; zwartość
przedziału zapewnia wspólną dziedzinę.
Wzór~\eqref{eq:druga-wariacja-energii}
daje (ii).

Dowiedźmy (iii) przez sprzeczność, pokazując
konkretny kierunek ujemnej drugiej wariacji.
Załóżmy, że $0<c<\ell$ i istnieje niezerowe
pole Jacobiego $J$ z $J(0)=J(c)=0$.
Przedłużmy je przez zero na $[c,\ell]$;
powstałe ciągłe pole oznaczmy $\widehat J$.
Jest gładkie oddzielnie po obu stronach $c$.
Z jednoznaczności równania Jacobiego wynika
$D_tJ(c)\ne0$, gdyż w przeciwnym razie
oba dane $J(c),D_tJ(c)$ byłyby zerowe.

Całkując przez części na $[0,c]$ i używając
równania Jacobiego, dla każdego pola $W$
o zerowych końcach otrzymujemy
\begin{equation}\label{eq:jacobi-przedluzone-indeks}
 I(\widehat J,W)
 =\bigl[g(D_tJ,W)\bigr]_0^c
 =g(D_tJ(c),W(c)),
 \qquad I(\widehat J,\widehat J)=0.
\end{equation}
Wybierzmy gładkie $W$ z $W(0)=W(\ell)=0$
i $W(c)=-D_tJ(c)$, na przykład przez
pomnożenie pola z lokalnej równoległej ramy
przez funkcję z nośnikiem blisko $c$.
Wtedy $I(\widehat J,W)=-|D_tJ(c)|^2<0$.
Dla dostatecznie małego $\varepsilon>0$
\[
 I(\widehat J+\varepsilon W,
   \widehat J+\varepsilon W)
 =-2\varepsilon|D_tJ(c)|^2
       +\varepsilon^2 I(W,W)<0.
\]
Wygładzamy $\widehat J$ w małym otoczeniu $c$
w równoległej ramie. Ponieważ jest ciągłe,
odcinkami gładkie i ma ograniczoną pochodną,
można to zrobić tak, by pole i jego pierwsza
pochodna zbliżały się w całce z kwadratu,
zachowując zerowe końce. Na zwartym przedziale
współczynniki krzywizny są ograniczone,
więc całka~\eqref{eq:forma-indeksowa}
zmienia się wtedy dowolnie mało. Otrzymujemy
\emph{gładkie} pole z $I<0$, sprzeczne z (ii).
Zatem punkt sprzężony nie może leżeć
we wnętrzu odcinka minimalnego.
Wreszcie przy $c=\ell$ wyraz $W(c)$ w
\eqref{eq:jacobi-przedluzone-indeks} musi
być zerowy, więc ten argument nie daje
ujemnego kierunku. Przykład sfery powyżej
pokazuje, że końcowy punkt sprzężony jest
możliwy.
\end{proof}

Brak punktów sprzężonych jest \emph{konieczny}
dla minimalności na otwartym odcinku,
ale sam nie zapewnia globalnie najkrótszej drogi.
Widać to nawet przy zerowej krzywiźnie.

\begin{przyklad}[Płaski walec: brak sprzężenia nie wystarcza]
\label{prz:jacobi-plaski-walec}
Na walcu $S_R^1\times\R$ z metryką
$g=R^2\dd\theta^2+\dd z^2$ poziomy łuk
$\gamma(t)=(t/R,0)$, $0\leq t\leq\ell$,
jest jednostkową geodezyjną. W lokalnych
współrzędnych metryka ma stałe współczynniki,
więc $R(X,Y)Z=0$; pola Jacobiego o zerowym
początku mają postać $tP_tw$ i dla $w\ne0$ nie znikają
ponownie. Gdy $\pi R<\ell<2\pi R$, drugi
łuk tego samego okręgu łączy te same końce
i ma długość $2\pi R-\ell<\ell$.
Geodezyjna przestaje zatem być globalnie
minimalna \emph{bez punktu sprzężonego}.
Przeszkodą jest tu globalne zawinięcie walca,
którego lokalne równanie Jacobiego nie widzi.
\end{przyklad}

\begin{figure}[htbp]
\centering
\begin{tikzpicture}[>=Latex,scale=1]
  \fill[manifoldblue!9] (0,0) circle (1.60);
  \draw[manifoldblue!75!black,very thick] (0,0) circle (1.60);
  \draw[manifoldred,ultra thick,->]
       (140:1.60) arc[start angle=140,end angle=340,radius=1.60];
  \draw[manifoldgreen!75!black,ultra thick,->]
       (140:1.60) arc[start angle=140,end angle=-20,radius=1.60];
  \fill[manifoldblue!85!black] (140:1.60) circle (2.5pt)
       node[above left] {$p$};
  \fill[manifoldblue!85!black] (-20:1.60) circle (2.5pt)
       node[below right] {$q$};
  \node[manifoldred!85!black] at (-2.75,-.55)
       {$\ell>\pi R$};
  \node[manifoldgreen!65!black] at (2.82,.57)
       {$2\pi R-\ell$};
  \node[manifoldblue!75!black] at (0,-2.08)
       {przekrój poziomy walca};
\end{tikzpicture}
\caption{Ten sam okrąg oglądany jako poziomy przekrój płaskiego walca. Między $p$ i $q$ dłuższy łuk (czerwony) jest geodezyjną bez punktów sprzężonych, lecz krótszy łuk (zielony) wygrywa globalnie.}
\label{fig:jacobi-walec-dwa-luki}
\end{figure}

Podsumowując logiczną kolejność:
\emph{minimalność długości} wymusza równanie
geodezyjnej przez pierwszą wariację, a następnie
nieujemność formy indeksowej przez drugą.
Pola Jacobiego opisują sąsiednie geodezyjne
i wskazują pierwszą przeszkodę w utrzymaniu
minimalności na coraz dłuższych odcinkach.
Na sferze przeszkodą jest ogniskowanie;
na walcu może nią być topologia, choć
ogniskowania w ogóle nie ma.

\section{Twierdzenia porównawcze}
\label{sec:twierdzenia-porownawcze}

Znamy już dokładne rozwiązania równania Jacobiego
przy stałej krzywiźnie. Co można powiedzieć, gdy
krzywizna zmienia się z punktu na punkt, ale znamy
jej górną albo dolną granicę? \emph{Porównanie}
z modelem o stałej krzywiźnie zamienia nierówność
dotyczącą $R$ w informację o geodezyjnych,
odległościach, objętości lub topologii.

Przez tę sekcję zakładamy $n=\dim M\geq2$
i metrykę riemannowską. Dla $\kappa\in\R$
oznaczmy
\begin{equation}\label{eq:porownanie-s-kappa}
 s_\kappa(r)=
 \begin{cases}
  \sin(\sqrt\kappa\,r)/\sqrt\kappa,&\kappa>0,\\
  r,&\kappa=0,\\
  \sinh(\sqrt{-\kappa}\,r)/\sqrt{-\kappa},&\kappa<0,
 \end{cases}
 \qquad c_\kappa(r)=\frac{s_\kappa'(r)}{s_\kappa(r)}.
\end{equation}
To funkcja z~\eqref{eq:jacobi-trzy-krzywizny}:
$s_\kappa''+\kappa s_\kappa=0$,
$s_\kappa(0)=0$, $s_\kappa'(0)=1$.
Ilorazu $c_\kappa$ używamy tylko tam, gdzie
$s_\kappa>0$; dla $\kappa>0$ oznacza to
$0<r<\pi/\sqrt\kappa$. Dla $\kappa=0$
mamy $c_0=1/r$, a dla $\kappa=-1$
mamy $c_{-1}=\coth r$.

\subsection*{Lemat indeksowy i porównanie Raucha}

Kluczowy krok jest wariacyjną własnością pola
Jacobiego. Piszemy $\mathsf K_t(X)=R(X,T)T$
na podprzestrzeni prostopadłej do $T=\dot\gamma$.
Ten operator jest symetryczny; jego forma
kwadratowa jest krzywizną sekcyjną płaszczyzny
rozpiętej przez $X,T$ razy $|X|^2$.

\begin{lemat}[Indeksowy]
\label{lem:indeksowy-porownanie}
Niech $\gamma:[0,b]\to M$ będzie jednostkową
geodezyjną bez punktu sprzężonego z $\gamma(0)$
na $(0,b]$. Niech $J$ będzie prostopadłym
polem Jacobiego z $J(0)=0$. Spośród wszystkich
ciągłych, odcinkami gładkich pól normalnych $X$ z
$X(0)=0$ i $X(b)=J(b)$ pole $J$ minimalizuje
formę indeksową na $[0,b]$:
\begin{equation}\label{eq:lemat-indeksowy}
 I_b(X,X)\geq I_b(J,J)
       =g(D_tJ(b),J(b)).
\end{equation}
\end{lemat}
\begin{proof}
Wybierzmy równoległą bazę ortonormalną
normalną do $T$ i zapisujmy pola przez
kolumny w $\R^{n-1}$. Niech $A(t)$ będzie
macierzą pól Jacobiego spełniających
$A(0)=0$ i $A'(0)=\mathrm{Id}$.
Każda kolumna spełnia
$A''+\mathsf K A=0$, gdzie $\mathsf K$
jest symetryczną macierzą. Brak punktów
sprzężonych oznacza, że $A(t)$ jest odwracalna
dla $0<t\leq b$: wektor z jądra $A(t)$
wyznaczałby niezerowe pole Jacobiego zerujące
się także w $t$.

Połóżmy $H=A'A^{-1}$. Najpierw sprawdźmy,
że $H$ jest symetryczna. Macierz Wrońskiego
$A^{\mathsf T}A'-A'^{\mathsf T}A$ ma pochodną
$A^{\mathsf T}A''-A''^{\mathsf T}A=0$,
bo $A''=-\mathsf K A$ i $\mathsf K^{\mathsf T}
=\mathsf K$. W zerze Wroński znika, więc
znika wszędzie. Mnożenie przez
$(A^{-1})^{\mathsf T}$ z lewej i $A^{-1}$
z prawej daje $H=H^{\mathsf T}$.
Różniczkując $H$ i korzystając z równania
Jacobiego, otrzymujemy \emph{równanie Riccatiego}
\begin{equation}\label{eq:jacobi-riccati}
 H'=A''A^{-1}-A'A^{-1}A'A^{-1}
    =-\mathsf K-H^2.
\end{equation}

Napiszmy $X=J+Z$; pole $Z$ znika na obu
końcach. Całkowanie przez części i równanie
Jacobiego dają $I_b(J,Z)=0$ oraz
$I_b(J,J)=[g(D_tJ,J)]_0^b$.
Z~\eqref{eq:jacobi-riccati} i symetrii $H$
wynika z kolei, dla $0<\varepsilon<b$,
\[
 \int_\varepsilon^b|Z'-HZ|^2\,\dd t
 =I_{[\varepsilon,b]}(Z,Z)
   -[\langle HZ,Z\rangle]_{\varepsilon}^{b}.
\]
Istotnie rozwijamy kwadrat i zastępujemy
$2\langle HZ,Z'\rangle$ wyrażeniem
$\frac{\dd}{\dd t}\langle HZ,Z\rangle-\langle H'Z,Z\rangle$;
pozostałe $H'+H^2$ to $-\mathsf K$.
Ponieważ $A(t)=t\mathrm{Id}+O(t^3)$,
to $H(t)=t^{-1}\mathrm{Id}+O(t)$.
Dla odcinkami gładkiego $Z$ z $Z(0)=0$
mamy $Z(t)=O(t)$, zatem wyraz brzegowy
przy $\varepsilon\downarrow0$ dąży do zera.
W rezultacie
\[
 I_b(Z,Z)=\int_0^b|Z'-HZ|^2\,\dd t\geq0.
\]
Razem z $I_b(X,X)=I_b(J,J)+I_b(Z,Z)$
daje to obie części~\eqref{eq:lemat-indeksowy}.
\end{proof}

Najpierw podajemy wygodną wersję Raucha:
porównujemy wszystkie płaszczyzny poprzeczne
obu geodezyjnych. Ta wersja wystarcza do poniższych
zastosowań; dowód używa właśnie podanego porównania
dla każdej pary kierunków poprzecznych.

\begin{twierdzenie}[Porównanie Raucha]
\label{thm:rauch-porownanie}
Niech $\gamma_i:[0,b)\to M_i$, $i=1,2$,
będą jednostkowymi geodezyjnymi w rozmaitościach
tego samego wymiaru. Załóżmy, że dla każdej
chwili $t$ i dowolnych jednostkowych wektorów
$X_i\perp\dot\gamma_i(t)$ zachodzi
\[
 K_1(\dot\gamma_1\wedge X_1)
 \leq K_2(\dot\gamma_2\wedge X_2).
\]
Niech normalne pola Jacobiego $J_i$ spełniają
$J_i(0)=0$ i
$|D_tJ_1(0)|=|D_tJ_2(0)|>0$.
Jeśli wzdłuż $\gamma_2$ nie ma punktu
sprzężonego z początkiem w $(0,b)$, to
\begin{equation}\label{eq:rauch-nierownosc}
 |J_1(t)|\geq |J_2(t)|,
 \qquad 0<t<b.
\end{equation}
W szczególności niższa krzywizna sekcyjna
pozwala geodezyjnym rozchodzić się co najmniej
tak szybko, jak w przestrzeni bardziej zakrzywionej.
\end{twierdzenie}
\begin{proof}
Przy ustalonym $t$ przed ewentualnym pierwszym
zerem $J_1$ połóżmy
$Y_i(s)=J_i(s)/|J_i(t)|$ dla $0\leq s\leq t$.
Wybierzmy izometrię między przestrzeniami
normalnymi obu geodezyjnych, przenoszoną
równolegle wzdłuż $s$, która przy $s=t$
odwzorowuje $Y_1(t)$ na $Y_2(t)$.
Niech $\widetilde Y_1$ będzie obrazem całego
pola $Y_1$ wzdłuż $\gamma_2$. Przenoszenie
zachowuje normy pól i ich pochodnych.
Założona nierówność krzywizn daje więc,
po odjęciu członów krzywizny w~\eqref{eq:forma-indeksowa},
\[
 I_{1,t}(Y_1,Y_1)
 \geq I_{2,t}(\widetilde Y_1,\widetilde Y_1).
\]
Pola $\widetilde Y_1$ oraz $Y_2$ mają te
same wartości: zero w $0$ i ten sam wektor
jednostkowy w $t$. Lemat~\ref{lem:indeksowy-porownanie}
dla $\gamma_2$ daje
$I_{2,t}(\widetilde Y_1,\widetilde Y_1)
\geq I_{2,t}(Y_2,Y_2)$.
Ponieważ $Y_i$ jest polem Jacobiego,
\eqref{eq:lemat-indeksowy} przekształca ten
łańcuch nierówności w
\[
 \frac{\dd}{\dd t}\log|J_1(t)|
 =\frac{g(D_tJ_1,J_1)}{|J_1|^2}
 \geq\frac{g(D_tJ_2,J_2)}{|J_2|^2}
 =\frac{\dd}{\dd t}\log|J_2(t)|.
\]
Całkujemy od $\varepsilon$ do $t$.
Przy $\varepsilon\downarrow0$ zachodzi
$|J_i(\varepsilon)|/\varepsilon
\to|D_tJ_i(0)|$, więc iloraz początkowy
dąży do $1$. Otrzymujemy~\eqref{eq:rauch-nierownosc}.
Jeżeli $J_1$ miałoby pierwsze zero przed $b$,
ta nierówność i ciągłość dawałyby tam
$0\geq|J_2|>0$, sprzeczność. Dowód obejmuje
zatem cały wskazany przedział.
\end{proof}

\begin{wniosek}[Porównanie z modelem]
\label{cor:rauch-model}
Jeżeli $K_M\leq\kappa$, to dla normalnego
pola Jacobiego $J(0)=0$, $D_tJ(0)=w$ mamy
\begin{equation}\label{eq:jacobi-gorne-K}
 |J(r)|\geq s_\kappa(r)|w|
\end{equation}
do pierwszego zera $s_\kappa$ (bez ograniczenia
czasu, gdy $\kappa\leq0$). Nie ma więc
punktu sprzężonego \emph{przed} tym zerem.
Jeżeli $K_M\geq\kappa$, nierówność jest
odwrotna aż do pierwszego punktu sprzężonego
w $M$ i pierwszego zera $s_\kappa$.
\end{wniosek}
\begin{proof}
Jako drugą przestrzeń w
Twierdzeniu~\ref{thm:rauch-porownanie}
weźmy model o stałej krzywiźnie $\kappa$.
Jego normalne pole ma długość
$s_\kappa(r)|w|$, zgodnie z
\eqref{eq:jacobi-trzy-krzywizny}.
Przy odwrotnym ograniczeniu krzywizny
zamieniamy role $M$ i modelu; przed
pierwszym punktem sprzężonym w $M$ spełniony
jest warunek lematu indeksowego.
\end{proof}

\begin{figure}[htbp]
\centering
\begin{tikzpicture}
\begin{axis}[
 width=10.7cm,height=6.0cm,
 axis lines=left,
 xmin=0,xmax=2.45,ymin=0,ymax=5.3,
 xtick={0,1,2},ytick={0,1,2,3,4,5},
 xlabel={$r$},ylabel={oddalenie $|J(r)|$},
 tick style={manifoldblue!65!black},
 label style={font=\small},
 clip=false]
 \addplot[manifoldred,very thick,domain=0:2.35,samples=90]
    {sin(deg(x))};
 \addplot[manifoldgold!80!black,very thick,domain=0:2.35]
    {x};
 \addplot[manifoldgreen!75!black,very thick,domain=0:2.35,samples=90]
    {(exp(x)-exp(-x))/2};
 \node[manifoldred!85!black,anchor=west]
    at (axis cs:2.08,.36) {$K=+1$};
 \node[manifoldgold!75!black,anchor=west]
    at (axis cs:2.08,1.80) {$K=0$};
 \node[manifoldgreen!65!black,anchor=west]
    at (axis cs:2.04,3.48) {$K=-1$};
\end{axis}
\end{tikzpicture}
\caption{W modelach o krzywiznach $+1,0,-1$ zerowe początkowe wychylenie i jednostkowa pochodna $D_rJ(0)$ dają odpowiednio $|J(r)|=\sin r$, $r$, $\sinh r$. Wykres pokazuje zakres $0\leq r\leq2{,}35$; pierwsze ponowne zero na sferze występuje później, przy $r=\pi$.}
\label{fig:rauch-porownanie-jacobiego}
\end{figure}

\subsection*{Porównanie hesjanu odległości}

W tej podsekcji zakładamy, że $(M,g)$ jest spójna
i zupełna. Najpierw uzasadnijmy, gdzie odległość
$r(x)=d_g(p,x)$ jest gładka.

\begin{lemat}[Czas cięcia i współrzędne biegunowe]
\label{lem:czas-ciecia-polar}
Dla $\theta\in S_pM$ połóżmy
\[
 c(\theta)=\sup\{t>0:d_g(p,\exp_p(t\theta))=t\}
 \in(0,\infty],\qquad
 C_p=\{\exp_p(c(\theta)\theta):c(\theta)<\infty\}.
\]
Zbiór $C_p$ nazywamy \emph{zbiorem cięcia} punktu $p$.
Odwzorowanie biegunowe $F(t,\theta)=\exp_p(t\theta)$
jest dyfeomorfizmem zbioru
$D=\{(t,\theta):0<t<c(\theta)\}$ na
$M\setminus(\{p\}\cup C_p)$. Na tym obrazie $r$ jest
gładka, a zbiór $C_p$ ma objętość riemannowską zero.
\end{lemat}
\begin{proof}
Małe promienie są minimalne, zatem $c(\theta)>0$.
Jeśli promień jest minimalny do chwili $t$, jest też
minimalny do każdej wcześniejszej chwili. Ciągłość
odległości zapewnia minimalność także przy skończonym
$c(\theta)$. Dla każdego $t>0$ zbiór
$\{\theta:c(\theta)\geq t\}
=\{\theta:d_g(p,F(t,\theta))=t\}$ jest domknięty,
więc $c$ jest mierzalna.

Przed czasem cięcia minimalny promień można przedłużyć
minimalnie nieco dalej. Dwie różne minimalne geodezyjne
do jego punktu końcowego utworzyłyby, po doklejeniu
tego przedłużenia do drugiej drogi, minimalną drogę
z narożnikiem. Jest to niemożliwe; gdyby prędkości
w miejscu sklejenia były równe, jednoznaczność ODE
dawałaby tę samą geodezyjną. Stąd jedyność, a także
różnowartościowość $F$ na $D$. Twierdzenie
\ref{thm:minimalna-bez-sprzezonego-wewnatrz} wyklucza
tam punkty sprzężone, więc $\dd\exp_p$ jest odwracalna.

Potrzebujemy jeszcze lokalnej kontroli minimalnych dróg.
Jeżeli do $q\ne p$ prowadzi jedyna minimalna geodezyjna
z wektorem początkowym $v$ i $\dd(\exp_p)_v$ jest
odwracalna, to dla $q_j\to q$ wektory minimalnych
geodezyjnych $v_j$ mają długości $d_g(p,q_j)\to d_g(p,q)$.
Istnieją na mocy Hopfa--Rinowa; zwartość ograniczonych
zbiorów w $T_pM$ i ciągłość $\exp_p$ pokazują, że każdy
zbieżny podciąg ma granicę $v$. Zatem wszystkie $v_j$
leżą ostatecznie w dziedzinie tej samej lokalnej
odwrotności $\exp_p$. W otoczeniu $q$ każda minimalna
droga jest więc wyznaczona przez tę odwrotność,
a $r(q')=|\exp_p^{-1}(q')|$ jest gładka.
W szczególności promienie dla danych bliskich $(t,\theta)$
z $D$ pozostają minimalne jeszcze przez krótki czas;
$D$ jest otwarty. Ten sam argument dowodzi, że punkt
o jedynej regularnej minimalnej geodezyjnej nie może
być końcem cięcia: promień dałoby się przedłużyć minimalnie.
Hopf--Rinow zapewnia teraz, że obraz $D$ jest dokładnie
podanym dopełnieniem, i lokalne odwrotności sklejają
się w globalną odwrotność $F$.

Pozostaje miara. Wykres skończonych wartości mierzalnej
funkcji $c$, czyli $G=\{(c(\theta),\theta):c(\theta)<\infty\}$,
ma miarę produktową zero: każdy jego przekrój w kierunku
$t$ jest jednopunktowy, więc stosujemy twierdzenie Fubiniego.
W mapach na sferze oraz na przedziałach
$1/k\leq t\leq k$ mapa $F$ jest lokalnie Lipschitz.
Takie odwzorowania między przestrzeniami tego samego
wymiaru przenoszą zbiory miary zero na zbiory miary zero
(pokrywamy zbiór małymi sześcianami, których obrazy
mają objętości ograniczone stałą razy objętości sześcianów).
Przeliczalne pokrycie mapami i porównanie gęstości
riemannowskiej z Lebesgue'a dają $\mu_g(F(G))=0$,
czyli $\mu_g(C_p)=0$.
\end{proof}

Poza $p$ i zbiorem cięcia lemat
Gaussa~\ref{lem:gaussa-exp} daje
$\nabla r=T$, gdzie $T$ jest jednostkowym
polem radialnych geodezyjnych. \emph{Hesjan}
funkcji $r$ jest formą dwuliniową
$\operatorname{Hess}r(X,Y)=g(\nabla_X\nabla r,Y)$.
Różni się od macierzy zwykłych drugich pochodnych
o symbole Christoffela i dlatego jest
niezależny od wyboru współrzędnych.

\begin{wniosek}[Porównanie hesjanu]
\label{cor:hess-odleglosc-porownanie}
W punktach, w których $r$ jest gładka,
jeżeli $K\leq\kappa$, to dla
$r<\pi/\sqrt\kappa$ przy $\kappa>0$
(i dla każdego $r>0$ przy $\kappa\leq0$)
\begin{equation}\label{eq:hess-porownanie}
 \operatorname{Hess}r(X,X)
 \geq c_\kappa(r)
       \bigl(|X|^2-g(X,\nabla r)^2\bigr).
\end{equation}
Jeżeli $K\geq\kappa$, nierówność się odwraca
przed zbiorem cięcia i pierwszym zerem
$s_\kappa$. Na modelu jest równość.
\end{wniosek}
\begin{proof}
Przypadek $X=0$ jest natychmiastowy. Weźmy najpierw
$0\ne X\perp T$ w końcu radialnej
geodezyjnej długości $r$. Ponieważ $\exp_p$
jest tu lokalnym dyfeomorfizmem, istnieje
normalne pole Jacobiego $J$ z $J(0)=0$
i $J(r)=X$; powstaje ono z wariacji
\emph{jednostkowych} prędkości początkowych.
Jej krzywe $t\mapsto F(s,t)$ są jednostkowymi
geodezyjnymi, więc dla ich pola stycznego
$T=\partial_tF$ mamy $D_sT=D_tJ$.
Różniczkując $\nabla r=T$ w kierunku
$X=J(r)$, otrzymujemy
\begin{equation}\label{eq:hess-jacobi}
 \operatorname{Hess}r(X,X)
 =g(D_tJ(r),J(r))
 =|X|^2\left.\frac{\dd}{\dd t}
               \log|J(t)|\right|_{t=r}.
\end{equation}
Rauch dowodził nawet nierówności między
pochodnymi logarytmów. Dla modelu
$|J_{\kappa}(t)|=s_\kappa(t)|w|$,
więc jego pochodna to $c_\kappa(t)$.
Daje to~\eqref{eq:hess-porownanie}
dla $X\perp T$ oraz nierówność odwrotną
po zamianie ról modeli. Ponieważ $|T|=1$,
$\operatorname{Hess}r(T,\cdot)=0$;
rozkładamy dowolne $X$ na część radialną
i prostopadłą i otrzymujemy pełną tezę.
\end{proof}

Na sferze $S_R^n$ prawa strona jest
$R^{-1}\cot(r/R)$, w $\R^n$ wynosi $1/r$,
a w $\mathbb H^n$ jest $\coth r$.
Na sferze zmienia znak za równikiem:
sfery geodezyjne przestają się rozszerzać
i kurczą się przed antypodą. W przestrzeni
hiperbolicznej rozszerzają się szybciej niż
w płaskiej, zgodnie z Rysunkiem~\ref{fig:rauch-porownanie-jacobiego}.

\subsection*{Dodatni Ricci: twierdzenie Bonneta--Myersa}

W poprzednich porównaniach ograniczaliśmy
krzywiznę \emph{każdej} płaszczyzny zawierającej
$T$. Poniższy wynik wymaga tylko dolnej granicy
\emph{sumy} tych krzywizn, czyli $\mathrm{Ric}(T,T)$
z~\eqref{eq:ricci-suma-sekcji}. Już ta suma
ogranicza, jak daleko może dojść minimalna
geodezyjna.

\begin{twierdzenie}[Bonneta--Myersa]
\label{thm:bonnet-myers}
Niech $(M^n,g)$ będzie spójna i zupełna, $n\geq2$.
Jeśli dla pewnego $\kappa>0$ i każdego
$v\in TM$ zachodzi
\begin{equation}\label{eq:myers-ricci-warunek}
 \mathrm{Ric}(v,v)\geq(n-1)\kappa|v|^2,
\end{equation}
to $M$ jest zwarta i
\begin{equation}\label{eq:myers-srednica}
 \operatorname{diam}(M,d_g)
 \leq\frac{\pi}{\sqrt\kappa}.
\end{equation}
Ponadto grupa podstawowa $\pi_1(M)$ jest skończona.
\end{twierdzenie}
\begin{proof}
Wybierzmy dowolne różne $p,q$ (dla $p=q$ oszacowanie
odległości jest oczywiste). Dzięki zupełności
Wniosek~\ref{cor:hopf-rinow-minimalna}
daje jednostkową geodezyjną minimalną
$\gamma:[0,\ell]\to M$ od $p$ do $q$.
Wybierzmy równoległe pola ortonormalne
$E_1,\ldots,E_{n-1}$ prostopadłe do $T$.
Dla każdego $i$ weźmy pole wariacyjne
$V_i(t)=f(t)E_i(t)$, gdzie
$f(t)=\sin(\pi t/\ell)$; obie jego wartości
brzegowe są zerowe. Z~\eqref{eq:forma-indeksowa}
otrzymujemy
\[
 I(V_i,V_i)=\int_0^\ell
 \left(f'(t)^2-f(t)^2g(R(E_i,T)T,E_i)\right)\dd t.
\]
Sumujemy po $i$. Suma wyrazów z krzywizną
to $\mathrm{Ric}(T,T)$, a
$\mathrm{Ric}(T,T)\geq(n-1)\kappa$, więc
\begin{align*}
 \sum_{i=1}^{n-1}I(V_i,V_i)
 &\leq(n-1)\int_0^\ell
     \bigl(f'(t)^2-\kappa f(t)^2\bigr)\dd t\\
 &=\frac{(n-1)\ell}{2}
           \left(\frac{\pi^2}{\ell^2}-\kappa\right).
\end{align*}
Użyliśmy tutaj
$\int_0^\ell\sin^2(\pi t/\ell)\dd t
=\int_0^\ell\cos^2(\pi t/\ell)\dd t=\ell/2$.
Gdyby $\ell>\pi/\sqrt\kappa$, prawa strona
byłaby ujemna. Co najmniej jeden składnik
$I(V_i,V_i)$ byłby wtedy ujemny, wbrew
Twierdzeniu~\ref{thm:minimalna-bez-sprzezonego-wewnatrz}.
Zatem $d_g(p,q)=\ell\leq\pi/\sqrt\kappa$.
Ponieważ $p,q$ były dowolne, mamy
\eqref{eq:myers-srednica}. Zbiór $M$ jest
domknięty i ograniczony jako przestrzeń
metryczna. Wersja Hopfa--Rinowa
z~\ref{thm:hopf-rinow-geodezyjne} mówi,
że na zupełnej rozmaitości riemannowskiej
takie zbiory są zwarte; stąd zwartość $M$.

Na koniec rozważmy uniwersalne nakrycie
$\widetilde M\to M$ z metryką podniesioną.
Jest ono zupełne: każda geodezyjna na nakryciu
rzutuje się na geodezyjną w $M$, którą można
przedłużyć na całą prostą, a przedłużenie
podnosi się jednoznacznie. Lokalna izometria
zachowuje Ricciego, więc ten sam dowód pokazuje,
że $\widetilde M$ jest zwarta. Włókno nakrycia
nad $p$ jest zamkniętym, dyskretnym podzbiorem
zwartej $\widetilde M$, a zatem jest skończone.
Liczba jego punktów równa się liczbie elementów
$\pi_1(M)$, co kończy dowód.
\end{proof}

\begin{przyklad}[Ostrość oszacowania]
Na $S_R^n$ mamy
$\mathrm{Ric}=(n-1)R^{-2}g$ i średnicę
$\pi R$; dla $\kappa=R^{-2}$ zachodzi
równość w~\eqref{eq:myers-srednica}.
Samo $\mathrm{Ric}>0$ punktowo nie daje
jednej liczby $\kappa$ na rozmaitości
niezwartej, jeśli nie ma \emph{jednostajnej}
dolnej granicy. Wymiar $n\geq2$ jest też
istotny: dla $n=1$ warunek
$\mathrm{Ric}\geq(n-1)\kappa g$ byłby pusty.
\end{przyklad}

\subsection*{Niedodatnia krzywizna: Cartan--Hadamard}

Ujemna krzywizna rozsuwa sąsiednie geodezyjne.
Gdy rozmaitość jest ponadto zupełna i nie ma
globalnego „zawinięcia” topologicznego,
odwzorowanie wykładnicze okazuje się
jedną mapą współrzędnych na całej rozmaitości.

\begin{twierdzenie}[Cartana--Hadamarda]
\label{thm:cartan-hadamard}
Jeżeli $(M,g)$ jest spójna, zupełna i ma
$K\leq0$, to dla każdego $p\in M$
odwzorowanie
$\exp_p:T_pM\to M$ jest nakryciem.
Jeśli ponadto $M$ jest jednospójna, to
$\exp_p$ jest globalnym dyfeomorfizmem;
w szczególności $M$ jest dyfeomorficzna
z $\R^n$ i każde dwa jej punkty łączy
jedyna geodezyjna minimalna.
\end{twierdzenie}
\begin{proof}
Zupełność i Hopf--Rinow dają
$D_p=T_pM$ oraz surjektywność $\exp_p$:
każdy $q$ leży na pewnej minimalnej
geodezyjnej startującej z $p$.
Z Wniosku~\ref{cor:rauch-model} dla $\kappa=0$
wynika, że niezerowa \emph{normalna} składowa
pola Jacobiego z $J(0)=0$ nie znika ponownie.
Składowa styczna ma postać $btT$, więc całe
niezerowe pole również nie może mieć drugiego zera.
Stwierdzenie~\ref{prop:jacobi-dexp} mówi,
że $\dd\exp_p$ nie ma jądra, więc $\exp_p$
jest lokalnym dyfeomorfizmem także daleko od $0$.

Pokażmy, dlaczego z tego wynika \emph{nakrycie},
a nie tylko lokalna odwracalność.
Wyposażmy $T_pM$ w metrykę $g_*:=\exp_p^*g$.
Jest to metryka riemannowska, a $\exp_p$
staje się lokalną izometrią. Niech $v=r\theta$
i rozłóżmy $w\in T_v(T_pM)$ na część radialną
$w_\parallel$ i prostopadłą $w_\perp$.
Lemat Gaussa pokazuje ortogonalność obrazów
tych części i zachowanie długości radialnej.
Dla poprzecznej części
$\dd(\exp_p)_v(w_\perp)$ jest wartością
przy $t=r$ pola Jacobiego o początkowej
pochodnej $w_\perp/r$. Rauch z porównaniem
do $\R^n$ daje
$|\dd(\exp_p)_v(w_\perp)|\geq|w_\perp|$.
Zatem, także w granicy dla $v=0$,
\begin{equation}\label{eq:hadamard-exp-rozszerza}
 g_{*,v}(w,w)=|\dd(\exp_p)_v(w)|_g^2
 \geq |w|_{g_p}^2.
\end{equation}
Dowolna droga w $(T_pM,g_*)$ ma długość
co najmniej równą długości euklidesowej.
Jej metryka odległości dominuje więc
euklidesową. Ciąg Cauchy'ego dla $g_*$
zbiega euklidesowo do pewnego $v$;
w małym otoczeniu $v$ obie gładkie metryki
są porównywalne, więc zbiega też dla $g_*$.
Dziedzina z metryką $g_*$ jest zupełna.

Przypomnijmy używany tutaj fakt o lokalnych
izometriach. Surjektywna lokalna izometria
$f:N\to M$ z \emph{zupełnego} spójnego $N$
jest nakryciem. Aby zobaczyć mechanizm,
weźmy małą normalną kulę $B_\rho(q)$ w $M$.
Dla każdego $x\in f^{-1}(q)$ radialne geodezyjne
z $q$ podnoszą się, dzięki zupełności $N$,
do radialnych geodezyjnych startujących w $x$.
Otrzymujemy gładką odwrotną gałąź nad
$B_\rho(q)$: pochodna $f_x$ odwraca
początkowy wektor, a wykładnicze odwzorowanie
w $N$ daje jego koniec. Ponieważ złożenie gałęzi z $f$
jest identycznością, jej różniczka jest odwracalna;
obraz gałęzi jest więc otwarty, a $f$ jest na nim
dyfeomorfizmem na kulę. Dwie takie gałęzie
nie mogą się przeciąć: cofnięcie podniesionej
geodezyjnej do $q$ wymusiłoby ten sam punkt $x$.
Każdy punkt przeciwobrazu kuli leży na którejś
gałęzi, bo można podnieść radialną drogę
od niego do $q$ i przedłużyć ją przez cały
krótki odcinek. Kula jest więc pokryta
rozłącznymi kopiami, co jest definicją nakrycia.
Stosujemy ten fakt do $f=\exp_p$.

Jeżeli $M$ jest jednospójna, spójne nakrycie
$T_pM\to M$ ma tylko jeden arkusz, zatem
jest dyfeomorfizmem. Hopf--Rinow dostarcza
geodezyjnej minimalnej między $p$ i każdym $q$;
różnowartościowość $\exp_p$ wymusza
jednoznaczność geodezyjnej od $p$ do $q$.
Ponieważ $p$ było dowolne, dotyczy to każdej
pary punktów.
\end{proof}

Hiperboloida oraz jej izometryczne modele
półprzestrzeni i dysku spełniają założenia
twierdzenia: dla nich $K=-1$ i $\exp_p$
jest globalnie odwracalna. Płaski walec
z Przykładu~\ref{prz:jacobi-plaski-walec}
pokazuje rolę jednospójności: jego $K=0$,
lecz $\exp_p$ jest nakryciem o wielu arkuszach,
a nie dyfeomorfizmem.

\subsection*{Porównanie objętości Bishopa--Gromowa}

Rauch porównuje pojedyncze kierunki, zaś
założenie o Riccim kontroluje ich \emph{sumę}.
W biegunowych współrzędnych wokół $p$
sumą tę odnajdziemy w pochodnej logarytmu
gęstości objętości.

\begin{twierdzenie}[Bishopa--Gromowa]
\label{thm:bishop-gromov}
Niech $(M^n,g)$ będzie spójna i zupełna, $n\geq2$,
oraz $\mathrm{Ric}\geq(n-1)\kappa g$.
Niech $V_p(r)=\mu_g(B_g(p,r))$ i
\begin{equation}\label{eq:objetosc-kappa-model}
 V_\kappa(r)
 =\omega_{n-1}\int_0^r s_\kappa(t)^{n-1}\dd t,
 \qquad \omega_{n-1}=\operatorname{vol}(S^{n-1}).
\end{equation}
W przypadku $\kappa>0$ rozważamy
$0<r<\pi/\sqrt\kappa$; w pozostałych
przypadkach dowolne $r>0$. Wówczas
\begin{equation}\label{eq:bishop-gromov}
 r\longmapsto\frac{V_p(r)}{V_\kappa(r)}
 \quad\text{jest nierosnąca},\qquad
 V_p(r)\leq V_\kappa(r).
\end{equation}
\end{twierdzenie}
\begin{proof}
\emph{Krok 1: gęstość wzdłuż jednego promienia.}
Dla jednostkowego $\theta\in T_pM$ piszemy
$\gamma_\theta(t)=\exp_p(t\theta)$.
Używamy czasu cięcia $c(\theta)$ z
Lematu~\ref{lem:czas-ciecia-polar}. Dla $0<t<c(\theta)$ nie ma punktu
sprzężonego na odcinku; macierz Jacobiego
$A(t)$ z dowodu Lematu~\ref{lem:indeksowy-porownanie}
jest odwracalna. Oznaczmy przez
$j(t,\theta)=\det A(t)>0$ biegunową gęstość
objętości. Jest dodatnia, bo
$A(t)=t\mathrm{Id}+O(t^3)$, i nie może zmienić
znaku bez wyzerowania wyznacznika.
Niech $H=A'A^{-1}$ i $m=\operatorname{tr}H$.
Wzór na pochodną wyznacznika daje
$m=\partial_t\log j$. Biorąc ślad
\eqref{eq:jacobi-riccati}, dostajemy
\[
 m'=-\operatorname{tr}(H^2)-\mathrm{Ric}(T,T)
 \leq-\frac{m^2}{n-1}-(n-1)\kappa.
\]
Pierwsza nierówność to Cauchy--Schwarz
dla wartości własnych symetrycznej macierzy $H$:
$\sum\lambda_i^2\geq(\sum\lambda_i)^2/(n-1)$.
W modelu stałej krzywizny
$m_\kappa=(n-1)s_\kappa'/s_\kappa$
spełnia równość w tym równaniu, bo
$s_\kappa''=-\kappa s_\kappa$.
Po odjęciu i oznaczeniu $u=m-m_\kappa$
otrzymujemy
\begin{equation}\label{eq:bg-riccati-roznica}
 u'\leq-\frac{m+m_\kappa}{n-1}\,u.
\end{equation}
Przy $t\downarrow0$ rozwinięcie
$A=t\mathrm{Id}+O(t^3)$ daje
$m=(n-1)/t+O(t)$, a
$m_\kappa=(n-1)/t+O(t)$, zatem $u=O(t)$.
Czynnik całkujący dla~\eqref{eq:bg-riccati-roznica}
zachowuje się jak $t^2$ przy zerze.
Mnożymy nierówność przez ten dodatni czynnik
i przechodzimy z dolną granicą całkowania
do zera: iloczyn czynnika przez $u$
ma granicę $0$. Wynika $u\leq0$, a więc
\[
 \partial_t\log\frac{j(t,\theta)}
               {s_\kappa(t)^{n-1}}
 =m-m_\kappa\leq0.
\]
Stosunek $h_\theta(t)=j(t,\theta)/
s_\kappa(t)^{n-1}$ maleje od wartości $1$.

\emph{Krok 2: dlaczego wolno całkować tylko do
czasu cięcia.} Dla $t\geq c(\theta)$
przyjmijmy $h_\theta(t)=0$. Nadal jest
nierosnąca. Na mocy Lematu~\ref{lem:czas-ciecia-polar} odwzorowanie
biegunowe jest dyfeomorfizmem przed czasem cięcia,
a pominięty zbiór ma miarę zero. Wzór na zmianę
zmiennych można więc stosować na tym otwartym zbiorze,
a następnie całkować po wszystkich kierunkach dzięki
twierdzeniu Tonellego (całkowana funkcja jest nieujemna).
Otrzymujemy polarowy wzór
\begin{equation}\label{eq:bg-polar-objetosc}
 V_p(r)=\int_{S_pM}\int_0^r
      h_\theta(t)s_\kappa(t)^{n-1}
                  \dd t\,\dd\theta.
\end{equation}

\emph{Krok 3: iloraz objętości.} Prawie wszędzie
pochodna $V_p'(r)$ istnieje. Z
\eqref{eq:bg-polar-objetosc} wynika, że
\[
 \frac{V_p'(r)}{V_\kappa'(r)}
 =\frac1{\omega_{n-1}}
          \int_{S_pM}h_\theta(r)\,\dd\theta
 =:a(r).
\]
Każda $h_\theta$ jest nierosnąca,
więc $a$ także. Ponadto
$V_p(r)=\int_0^r a(t)V_\kappa'(t)\dd t$;
iloraz $V_p(r)/V_\kappa(r)$ jest zatem
średnią wartości $a(t)$ z dodatnią wagą
$V_\kappa'(t)$. Dla nierosnącej $a$
ta średnia jest co najmniej $a(r)$.
Reguła ilorazu daje prawie wszędzie
\[
 \left(\frac{V_p}{V_\kappa}\right)'
 =\frac{V_\kappa'}{V_\kappa}
   \left(a(r)-\frac{V_p}{V_\kappa}\right)\leq0.
\]
Ponieważ $0\leq a\leq1$, wzór całkowy pokazuje, że
$V_p$ jest lokalnie absolutnie ciągła. Na każdym zwartym
przedziale promieni dodatnich $V_\kappa$ jest gładka
i dodatnia, więc iloraz również jest absolutnie ciągły.
Całkując jego niedodatnią pochodną, otrzymujemy
monotoniczność wszędzie.
Sama ciągłość i nierówność dla pochodnej prawie wszędzie
nie wystarczyłyby do tego wniosku. Przy $r\downarrow0$ obie objętości
mają asymptotykę $\omega_{n-1}r^n/n$,
stąd iloraz dąży do $1$ i otrzymujemy
drugą część~\eqref{eq:bishop-gromov}.
\end{proof}

Ten sam rachunek daje \emph{porównanie laplasjanu
odległości}. Przy konwencji
$\Delta r=\operatorname{div}(\nabla r)
=\operatorname{tr}_g\operatorname{Hess}r$,
poza $p$ i zbiorem cięcia,
\begin{equation}\label{eq:laplasjan-porownanie}
 \mathrm{Ric}\geq(n-1)\kappa g
 \quad\Longrightarrow\quad
 \Delta r=\operatorname{tr}(A'A^{-1})
       =m\leq(n-1)c_\kappa(r).
\end{equation}
To właśnie nierówność z Kroku~1 dowodu
Bishopa--Gromowa. Gdy $K\leq\kappa$, ślad
\eqref{eq:hess-porownanie} daje odwrotną
nierówność $\Delta r\geq(n-1)c_\kappa(r)$.
Obie działają w odpowiednim zakresie
$s_\kappa(r)>0$ przed zbiorem cięcia.

\begin{przyklad}[Co widać w modelach]
Dla $\R^n$ przy $\kappa=0$ oraz
$\mathbb H^n$ przy $\kappa=-1$ iloraz
$V_p(r)/V_\kappa(r)$ jest stale równy $1$.
Na $S_R^n$ zachodzi to samo dla
$\kappa=R^{-2}$ i $r<\pi R$.
Na płaskim walcu Ricci jest zerowy, ale
objętość dużej kuli rośnie tylko liniowo
z promieniem, podczas gdy dla $\R^2$
rośnie jak $r^2$: sam warunek $\mathrm{Ric}\geq0$
nie wymusza euklidesowego \emph{dolnego}
tempa wzrostu objętości.
\end{przyklad}

\begin{uwaga}[Która krzywizna kontroluje co?]
Rauch i porównanie hesjanu korzystają z granic
na każdą \emph{krzywiznę sekcyjną}; Bonnet--Myers
i Bishop--Gromov tylko z granic na \emph{Ricciego},
czyli sumę takich krzywizn. Cartan--Hadamard
dodaje do $K\leq0$ zupełność i jednospójność,
by otrzymać globalną informację o topologii
i jednoznaczności geodezyjnych. Zamiana górnej
granicy na dolną odwraca kierunek porównania.
\end{uwaga}

\section{Twierdzenie Gaussa--Bonneta}
\label{sec:gauss-bonnet}
\index{twierdzenie Gaussa--Bonneta}

Dotąd krzywizna opisywała geometrię w pobliżu punktu:
jak rozchodzą się geodezyjne i jak zmienia się transport
równoległy. Teraz zsumujemy krzywiznę \emph{na całej
powierzchni}. Wynik zależy od liczby składowych, uchwytów
i brzegów powierzchni, czyli od jej topologii. Jest to
szczególnie wyraźne po porównaniu sfery z torusem.

W całej sekcji $S$ jest powierzchnią z \emph{metryką
riemannowską}, $K(p)=K_p(T_pS)$ jej krzywizną sekcyjną,
a $\dd A$ dodatnim elementem
pola z Sekcji~\ref{sec:metryka-objetosc}. Dla powierzchni
zanurzonej w $\R^3$ twierdzenie Gaussa utożsamia tę
krzywiznę z iloczynem krzywizn głównych.
Całkę $\int_S K\,\dd A$ można określić również wtedy,
gdy $S$ nie jest orientowalna: $K\,\dd A$ jest gęstością,
nie wymaga więc wyboru orientacji. W dowodzie najpierw
wybierzemy orientację.

\begin{definicja}[Krzywizna geodezyjna i skręt narożnika]
\index{krzywizna geodezyjna i skret naroznika@Krzywizna geodezyjna i skręt narożnika}
\label{def:krzywizna-geodezyjna}
Niech $S$ będzie zorientowana, a $J:T_pS\to T_pS$
oznacza dodatni obrót o $90^\circ$: jeżeli $(e_1,e_2)$
jest dodatnią bazą ortonormalną, to $Je_1=e_2$.
Dla łuku brzegu parametryzowanego długością $s$ oznaczmy
jednostkowy wektor styczny przez $T=\dot\gamma(s)$.
\emph{Krzywizna geodezyjna ze znakiem} tego łuku wynosi
\begin{equation}\label{eq:kg-def}
 k_g=g(\nabla_TT,JT).
\end{equation}
W punkcie narożnym niech $\beta\in(0,2\pi)$ oznacza
kąt \emph{wewnątrz} obszaru. Jego \emph{zewnętrzny skręt}
to $\varepsilon=\pi-\beta$; dla wklęsłego narożnika
jest ujemny. Brzeg obszaru orientujemy tak, aby wnętrze
znajdowało się po lewej stronie ruchu. Na łuku gładkim
$JT$ wskazuje wtedy lokalnie wnętrze obszaru.
\end{definicja}

Skąd w definicji znak? Na brzegu euklidesowego dysku
promienia $r$, przebieganym przeciwnie do ruchu wskazówek
zegara, $\nabla_TT$ wskazuje środek i ma długość $1/r$.
Wektor $JT$ także wskazuje środek, więc $k_g=1/r$.
Na brzegu otworu dodatni kierunek obiegu jest odwrotny:
$JT$ wskazuje wówczas na zewnątrz otworu, a $k_g$
okręgu jest ujemne. Równość $k_g=0$ znaczy, że łuk
jest geodezyjną (zob.\ Sekcję~\ref{sec:geodezyjna-pierwsza-wariacja}).

\begin{figure}[htbp]
\centering
\begin{tikzpicture}[>=Latex,font=\small]
 \begin{scope}[shift={(-3,-.8)}]
  \fill[manifoldblue!12] (0,0)--(2,0) arc[start angle=0,end angle=90,radius=2]--cycle;
  \draw[manifoldblue!78!black,thick,postaction={decorate,decoration={markings,
    mark=at position .12 with {\arrow{Latex}},
    mark=at position .87 with {\arrow{Latex}}}}]
   (0,0)--(2,0) arc[start angle=0,end angle=90,radius=2]--cycle;
  \coordinate (P) at ({sqrt(2)},{sqrt(2)});
  \fill[manifoldred] (P) circle (1.8pt);
  \draw[->,manifoldred,thick] (P)--({sqrt(2)-.55},{sqrt(2)+.55})
   node[above] {$T$};
  \draw[->,manifoldgreen!75!black,thick] (P)--({sqrt(2)-.55},{sqrt(2)-.55})
   node[right] {$JT$};
  \draw[manifoldgold!75!black] (.35,0) arc[start angle=0,end angle=90,radius=.35];
  \node at (.48,.47) {$\beta$};
  \node[manifoldblue!75!black] at (.72,1.22) {$\Omega$};
  \node[below] at (1,-.45) {dodatnio zorientowany brzeg};
 \end{scope}
 \begin{scope}[shift={(2.6,.1)}]
  \shade[inner color=white,outer color=manifoldblue!43] (0,0) circle (1.48);
  \draw[manifoldblue!76!black,thick] (0,0) circle (1.48);
  \fill[manifoldgold!35] (0,0) circle ({1.48*sin(60)});
  \draw[manifoldgreen!70!black,very thick,postaction={decorate,
   decoration={markings,mark=at position .18 with {\arrow{Latex}},
   mark=at position .68 with {\arrow{Latex}}}}]
   (0,0) circle ({1.48*sin(60)});
  \fill[manifoldblue!80!black] (0,0) circle (1.6pt) node[below] {biegun};
  \node at (0,.56) {$\Omega_{\pi/3}$};
  \node[below] at (0,-1.65) {widok od bieguna północnego};
 \end{scope}
\end{tikzpicture}
\caption{Po lewej ćwierćdysk z dodatnio zorientowanym brzegiem:
$T$ jest styczny do okręgu, a $JT$ skierowany do obszaru;
zaznaczony kąt wewnętrzny to $\beta=\pi/2$.
Po prawej rzut czapki sferycznej o kącie biegunowym $\theta_0=\pi/3$.
Jej brzeg obiegamy przeciwnie do ruchu wskazówek zegara,
zgodnie z orientacją od zewnętrznej normalnej sfery.}
\label{fig:gauss-bonnet-brzeg}
\end{figure}

\begin{twierdzenie}[Gaussa--Bonneta z brzegiem i narożnikami]
\label{thm:gauss-bonnet-brzeg}
Niech $\Omega$ będzie zwartą, zorientowaną powierzchnią
riemannowską z brzegiem złożonym ze skończonej liczby
regularnych łuków klasy $C^2$. Załóżmy, że w narożnikach
istnieją jednostronne styczne i kąty wewnętrzne
$\beta_1,\ldots,\beta_m\in(0,2\pi)$. Wówczas
\begin{equation}\label{eq:gauss-bonnet-brzeg}
 \boxed{\displaystyle
 \int_\Omega K\,\dd A
 +\int_{\partial\Omega}k_g\,\dd s
 +\sum_{j=1}^m(\pi-\beta_j)
 =2\pi\chi(\Omega).}
\end{equation}
Brzeg całkujemy odcinek po odcinku. Dla brzegu gładkiego
suma po narożnikach jest pusta. Charakterystykę Eulera
można obliczyć z triangulacji:
$\chi(\Omega)=V-E+F$, gdzie $V,E,F$ są liczbami
wierzchołków, krawędzi i trójkątów; zgadza się to z
opisem homologicznym z
Twierdzenia~\ref{thm:euler-homologia}.
\end{twierdzenie}

\begin{proof}
\emph{Krok 1: forma koneksji na małym obszarze.}
Załóżmy najpierw, że $D\subset\Omega$ jest topologicznym
dyskiem, leży w jednej dodatnio zorientowanej mapie i ma
brzeg odcinkami gładki. Na otoczeniu $D$ wybierzmy
dodatnią ortonormalną ramę $(e_1,e_2)$, $e_2=Je_1$.
Istnieje ona na małym obszarze współrzędnych: wektor
pierwszej osi można unormować, a drugi otrzymać przez $J$.
Definiujemy lokalną $1$-formę
\[
 \alpha(X)=g(\nabla_Xe_1,e_2).
\]
Zgodność koneksji z metryką, czyli
$Xg(e_i,e_j)=g(\nabla_Xe_i,e_j)+g(e_i,\nabla_Xe_j)$,
daje dokładnie
\begin{equation}\label{eq:gb-rama}
 \nabla_Xe_1=\alpha(X)e_2,
 \qquad \nabla_Xe_2=-\alpha(X)e_1.
\end{equation}
Obliczmy z tego tensor krzywizny przy przyjętym tu
znaku~\eqref{eq:riemann-def}. Współczynnik przy $e_2$
wektora $\nabla_X\nabla_Ye_1$ wynosi $X(\alpha(Y))$;
wyrazy przy $e_1$ z dwóch kolejności odejmują się.
Dlatego
\[
 R(X,Y)e_1
 =\bigl(X\alpha(Y)-Y\alpha(X)-\alpha([X,Y])\bigr)e_2
 =\dd\alpha(X,Y)e_2.
\]
Dla dodatniej ramy ortonormalnej definicja krzywizny
sekcyjnej mówi, że
$g(R(e_1,e_2)e_2,e_1)=K$. Ponieważ $R(e_1,e_2)$
jest antysymetryczny względem $g$, mamy również
$g(R(e_1,e_2)e_1,e_2)=-K$. Porównanie z poprzednią
równością daje
\begin{equation}\label{eq:gb-dalpha}
 \dd\alpha=-K\,\dd A.
\end{equation}
Minus w tym miejscu jest konsekwencją naszej konwencji
dla $R$; jego zgubienie odwróciłoby znak całki z $K$.

\emph{Krok 2: krzywizna łuku jako obrót stycznej.}
Na gładkim łuku brzegu dysku można wybrać ciągłą
funkcję kąta $\vartheta(s)$ taką, że
$T=\cos\vartheta\,e_1+\sin\vartheta\,e_2$.
Różniczkując ten wzór wzdłuż łuku i podstawiając
\eqref{eq:gb-rama}, otrzymujemy
\[
 \nabla_TT
 =\bigl(\vartheta'(s)+\alpha(T)\bigr)
             (-\sin\vartheta\,e_1+\cos\vartheta\,e_2)
 =\bigl(\vartheta'(s)+\alpha(T)\bigr)JT.
\]
Stąd $k_g\,\dd s=\dd\vartheta+\alpha$ na każdym łuku.
Przy przejściu przez narożnik kierunek stycznej
obraca się dodatkowo o $\pi-\beta_j$: w wypukłym
narożniku w lewo, we wklęsłym w prawo.

Suma zmian kąta $\vartheta$ na łukach wraz z obrotami
w narożnikach wynosi $2\pi$. Jest to twierdzenie o
obrocie stycznej do dodatnio zorientowanej krzywej
Jordana na płaszczyźnie. Żeby zobaczyć, skąd bierze się
liczba $2\pi$, można łuki przybliżyć łamanymi wraz z
ich kierunkami stycznymi. W prostym wielokącie o $v$
wierzchołkach można poprowadzić $v-3$ przekątnych
i otrzymać $v-2$ trójkąty. Suma kątów wewnętrznych
to $(v-2)\pi$, a suma skrętów zewnętrznych wynosi
$v\pi-(v-2)\pi=2\pi$. Przy zagęszczaniu łamanej skręty
wewnątrz łuków dążą do całek z $\vartheta'\dd s$;
skręty w pierwotnych narożnikach zostają. Wybór
dodatniej ramki nad dyskiem nie zmienia wyniku. Przejście
od ramki współrzędnych do ramki ortonormalnej jest mapą
do $\mathrm{GL}^+(2,\R)$, a nie samym obrotem. W rozkładzie
polarnym tej macierzy dodatnio określony czynnik deformujemy
do identyczności przez $(1-s)P+s\mathrm{Id}$; niezerowa
styczna pozostaje niezerowa, więc jej liczba obrotów się
nie zmienia. Pozostaje mapa obrotów do $\mathrm{SO}(2)$.
Przedłuża się ona na cały dysk, zatem na brzegu ma stopień
zero. Liczba obrotów stycznej w naszej ramce nadal wynosi $1$.

Sumujemy równość $k_g\dd s=\dd\vartheta+\alpha$
po łukach, dodajemy obroty narożników i stosujemy
Twierdzenie Stokesa~\ref{thm:stokes-formy} oraz
\eqref{eq:gb-dalpha}:
\begin{equation}\label{eq:gb-dysk}
 \int_{\partial D}k_g\,\dd s
 +\sum_j(\pi-\beta_j)
 =2\pi+\int_{\partial D}\alpha
 =2\pi-\int_D K\,\dd A.
\end{equation}
Stokesa dla brzegu z narożnikami można tu uzyskać
z wersji dla brzegu gładkiego przez zaokrąglenie
narożników: długość usuniętych łuków i pole usuniętych
małych fragmentów dążą do zera, a $\alpha$ i $\dd\alpha$
są ograniczone na zwartym otoczeniu. Zaokrąglamy
\emph{tylko} przy zastosowaniu Stokesa; powyższe
skręty narożników pozostają w rachunku kątów.
Tak otrzymaliśmy wzór dla dysku, ponieważ $\chi(D)=1$.

\emph{Krok 3: sklejenie lokalnych wzorów.}
Wybierzmy skończoną gładką triangulację $\Omega$ na tyle
drobną, aby każdy jej trójkąt leżał w jednej mapie;
narożniki brzegu włączmy do zbioru wierzchołków.
Można to zrobić, korzystając z triangulowalności
gładkich powierzchni, z regularnymi krawędziami i
niezdegenerowanymi kątami trójkątów przy wierzchołkach. Zastosujmy~\eqref{eq:gb-dysk}
do każdego z $F$ trójkątów i zsumujmy. Całki $K$
sumują się do całki po $\Omega$. Na wewnętrznej
krawędzi dwie orientacje są przeciwne, a wraz z nimi
zmienia się znak $k_g$, więc jej wkłady znoszą się.
Zostaje całka po rzeczywistym brzegu.

Oznaczmy liczbę wierzchołków wewnętrznych przez
$V_{\mathrm w}$, a brzegowych przez $V_{\mathrm b}$.
Przy każdym wewnętrznym wierzchołku suma kątów
trójkątów wynosi $2\pi$, a przy każdym brzegowym
wynosi kąt $\beta_v$ obszaru (jeśli brzeg jest tam
gładki, $\beta_v=\pi$). Suma zewnętrznych skrętów
\emph{wszystkich trójkątów} jest więc
$3\pi F-2\pi V_{\mathrm w}-\sum_{v\in\partial\Omega}\beta_v$.
Po zsumowaniu~\eqref{eq:gb-dysk} dostajemy
\begin{equation}\label{eq:gb-liczenie}
 \int_\Omega K\,\dd A
 +\int_{\partial\Omega}k_g\,\dd s
 =-\pi F+2\pi V_{\mathrm w}
                 +\sum_{v\in\partial\Omega}\beta_v.
\end{equation}
Każda krawędź wewnętrzna należy do dwóch trójkątów,
a brzegowa do jednego. Każda składowa brzegu jest
zamkniętym cyklem z taką samą liczbą krawędzi jak
wierzchołków. Zatem
$3F=2E-V_{\mathrm b}$ (w przypadku pustego brzegu
oba ostatnie zliczenia brzegowe są zerowe).
Ponieważ $V=V_{\mathrm w}+V_{\mathrm b}$,
\[
 2\pi(V-E+F)
 -\sum_{v\in\partial\Omega}(\pi-\beta_v)
 =-\pi F+2\pi V_{\mathrm w}
                  +\sum_{v\in\partial\Omega}\beta_v.
\]
Wstawiamy tę tożsamość do~\eqref{eq:gb-liczenie}.
Sztucznie wstawione wierzchołki na gładkich częściach
brzegu mają $\beta_v=\pi$ i niczego nie dodają.
Otrzymujemy~\eqref{eq:gauss-bonnet-brzeg}.
\end{proof}

\begin{wniosek}[Gauss--Bonnet dla powierzchni zamkniętej]
\label{cor:gauss-bonnet-zamknieta}
Jeżeli $S$ jest zwartą powierzchnią bez brzegu, to
\begin{equation}\label{eq:gb-zamknieta}
 \boxed{\displaystyle\int_S K\,\dd A=2\pi\chi(S).}
\end{equation}
W szczególności dla spójnej orientowalnej powierzchni
rodzaju $g$ prawa strona wynosi $4\pi(1-g)$.
\end{wniosek}
\begin{proof}
Dla powierzchni orientowalnej stosujemy
Twierdzenie~\ref{thm:gauss-bonnet-brzeg} przy
$\partial S=\varnothing$. Z klasyfikacji powierzchni
zamkniętych $\chi(S)=2-2g$; równoważnie można tę
liczbę policzyć z dowolnej triangulacji standardowego
modelu sfery z $g$ uchwytami. Jeśli $S$ nie jest
orientowalna, bierzemy jej dwukrotne nakrycie
orientacyjne $\widetilde S\to S$. Metryka podniesiona
ma te same lokalne wartości $K$, a każdy mały obszar
na $S$ występuje dwa razy na $\widetilde S$, więc
$\int_{\widetilde S}K\dd A=2\int_SK\dd A$.
Podnosząc triangulację, otrzymujemy też
$\chi(\widetilde S)=2\chi(S)$. Dzielimy zorientowany
wzór przez $2$.
\end{proof}

\begin{przyklad}[Dysk i czapka sferyczna]
\label{prz:gb-dysk-czapka}
Dysk euklidesowy $D_r$ ma $K=0$, $k_g=1/r$ i brzeg
długości $2\pi r$. Stąd
$\int_{\partial D_r} k_g\dd s=2\pi=2\pi\chi(D_r)$.
Gdyby pominąć całkę po brzegu, wyszłaby fałszywa
równość $0=2\pi$. Dla płaskiego pierścienia o
promieniach $0<r<R$ zewnętrzny okrąg wnosi $2\pi$,
wewnętrzny $-2\pi$: suma wynosi $0$, zgodnie z
$\chi(S^1\times[0,1])=0$.

Na sferze $S_R^2$ opiszmy czapkę od bieguna północnego
do kąta biegunowego (kolatytudy) $\theta_0\in(0,\pi)$:
\[
 0\leq\theta\leq\theta_0,\qquad
 g=R^2(\dd\theta^2+\sin^2\theta\,\dd\varphi^2),
 \qquad K=R^{-2}.
\]
Z elementu pola $R^2\sin\theta\,\dd\theta\dd\varphi$
dostajemy
\[
 \int_{\Omega_{\theta_0}}K\dd A
 =\int_0^{2\pi}\int_0^{\theta_0}\sin\theta
                                     \,\dd\theta\dd\varphi
 =2\pi(1-\cos\theta_0).
\]
Brzeg ma długość $2\pi R\sin\theta_0$. Jego dodatnio
skierowany styczny to $T=(R\sin\theta_0)^{-1}
\partial_\varphi$, a $JT=-R^{-1}\partial_\theta$.
Ponieważ $\Gamma^\theta_{\varphi\varphi}
=-\sin\theta\cos\theta$, wzór~\eqref{eq:kg-def} daje
\[
 k_g=\frac{\cos\theta_0}{R\sin\theta_0},
 \qquad \int_{\partial\Omega_{\theta_0}}k_g\dd s
 =2\pi\cos\theta_0.
\]
Suma obu całek to $2\pi$, zgodnie z $\chi=1$.
Dla półsfery $\theta_0=\pi/2$ równik jest geodezyjną
i cały wkład pochodzi z $K$; dla czapki większej
niż półsfera wkład brzegu jest \emph{ujemny}.
\end{przyklad}

\begin{przyklad}[Kąty trójkąta geodezyjnego]
\label{prz:gb-trojkat}
Niech trzy łuki geodezyjne tworzą prostą zamkniętą krzywą
ograniczającą obszar $\Delta$ homeomorficzny z domkniętym
dyskiem, o kątach wewnętrznych $A,B,C$.
Wtedy $k_g=0$ na każdym boku, a $\chi(\Delta)=1$.
Z~\eqref{eq:gauss-bonnet-brzeg} wynika
\begin{equation}\label{eq:gb-trojkat}
 \int_\Delta K\,\dd A=A+B+C-\pi.
\end{equation}
W płaszczyźnie daje to $A+B+C=\pi$. Na sferze
$S_R^2$ pole trójkąta to
$R^2(A+B+C-\pi)$. Trójkąt o wierzchołkach
na biegunie północnym i w dwóch punktach równika
oddalonych o ćwierć obwodu ma $A=B=C=\pi/2$;
jako jedna ósma sfery ma pole $\pi R^2/2$.
W płaszczyźnie hiperbolicznej $K=-1$,
\[
 \operatorname{Area}(\Delta)=\pi-(A+B+C)>0.
\]
To wprost pokazuje, jak dodatnia i ujemna krzywizna
zmieniają sumę kątów.
Rysunek~\ref{fig:gb-trojkaty-modele} zestawia
odpowiednie geodezyjne w trzech modelach.
\end{przyklad}

\begin{figure}[H]
\centering
\begin{tikzpicture}[>=Latex,line cap=round,line join=round,
                     font=\small]
  % Trójkąt euklidesowy, o dokładnie równych bokach.
  \begin{scope}[shift={(-4.05,0)}]
   \fill[manifoldblue!13]
       (-1.05,-.75)--(1.05,-.75)--(0,1.06865)--cycle;
   \draw[manifoldblue!85!black,very thick]
       (-1.05,-.75)--(1.05,-.75)--(0,1.06865)--cycle;
   \foreach \x/\y in {-1.05/-.75,1.05/-.75,0/1.06865}
      \fill[manifoldred] (\x,\y) circle (1.7pt);
   \node[manifoldblue!80!black] at (0,-1.40) {$\R^2,\ K=0$};
   \node at (0,-1.74) {$A+B+C=\pi$};
  \end{scope}
  % Oktant sferyczny widziany w rzucie prostopadłym z (1,1,1).
  \begin{scope}
   \shade[inner color=white,outer color=manifoldblue!43] (0,0) circle (1.25);
   \draw[manifoldblue!78!black,thick] (0,0) circle (1.25);
   \path[fill=manifoldgold!43,opacity=.83]
     plot[domain=0:90,samples=30]
       ({1.25*(cos(\x)-sin(\x))/sqrt(2)},
        {1.25*(cos(\x)+sin(\x))/sqrt(6)}) --
     plot[domain=0:90,samples=30]
       ({-1.25*cos(\x)/sqrt(2)},
        {1.25*(cos(\x)-2*sin(\x))/sqrt(6)}) --
     plot[domain=0:90,samples=30]
       ({1.25*sin(\x)/sqrt(2)},
        {1.25*(-2*cos(\x)+sin(\x))/sqrt(6)}) -- cycle;
   \draw[manifoldgold!75!black,very thick]
     plot[domain=0:90,samples=30]
       ({1.25*(cos(\x)-sin(\x))/sqrt(2)},
        {1.25*(cos(\x)+sin(\x))/sqrt(6)});
   \draw[manifoldgold!75!black,very thick]
     plot[domain=0:90,samples=30]
       ({-1.25*cos(\x)/sqrt(2)},
        {1.25*(cos(\x)-2*sin(\x))/sqrt(6)});
   \draw[manifoldgold!75!black,very thick]
     plot[domain=0:90,samples=30]
       ({1.25*sin(\x)/sqrt(2)},
        {1.25*(-2*cos(\x)+sin(\x))/sqrt(6)});
   \foreach \x/\y in {.88388/.51031,-.88388/.51031,0/-1.02062}
      \fill[manifoldred] (\x,\y) circle (1.7pt);
   \node[manifoldblue!80!black] at (0,-1.40) {$S_R^2,\ K=R^{-2}$};
   \node at (0,-1.74) {$A=B=C=\pi/2$};
  \end{scope}
  % Poincaré: boki są łukami okręgów prostopadłych do brzegu.
  \begin{scope}[shift={(4.05,0)}]
   \pgfmathsetmacro{\hc}{(1+.78^2)/.78}
   \pgfmathsetmacro{\hr}{sqrt(\hc^2-1)}
   \pgfmathsetmacro{\ha}{asin(.78*sin(60)/\hr)}
   \fill[manifoldgreen!5] (0,0) circle (1.25);
   \draw[manifoldgreen!60!black,densely dashed,thick]
        (0,0) circle (1.25);
   \path[fill=manifoldgreen!25]
    plot[domain=\ha:-\ha,samples=30]
      ({1.25*(\hc*cos(150)+\hr*cos(330+\x))},
       {1.25*(\hc*sin(150)+\hr*sin(330+\x))}) --
    plot[domain=\ha:-\ha,samples=30]
      ({1.25*(\hc*cos(270)+\hr*cos(450+\x))},
       {1.25*(\hc*sin(270)+\hr*sin(450+\x))}) --
    plot[domain=\ha:-\ha,samples=30]
      ({1.25*(\hc*cos(390)+\hr*cos(570+\x))},
       {1.25*(\hc*sin(390)+\hr*sin(570+\x))}) -- cycle;
   \foreach \m in {150,270,390}
     \draw[manifoldgreen!72!black,very thick]
       plot[domain=\ha:-\ha,samples=35]
       ({1.25*(\hc*cos(\m)+\hr*cos(\m+180+\x))},
        {1.25*(\hc*sin(\m)+\hr*sin(\m+180+\x))});
   \foreach \ang in {90,210,330}
     \fill[manifoldred] (\ang:.975) circle (1.7pt);
   \node[manifoldgreen!70!black] at (0,-1.40) {$\mathbb H^2,\ K=-1$};
   \node at (0,-1.74) {$A+B+C<\pi$};
  \end{scope}
\end{tikzpicture}
\caption{Trójkąty geodezyjne w geometrii euklidesowej,
na sferze i w modelu dyskowym płaszczyzny hiperbolicznej.
Środkowy trójkąt jest oktantem sfery: każdy jego bok
stanowi ćwiartkę wielkiego koła, a każdy kąt jest prosty.
W dysku boki leżą na okręgach prostopadłych do
przerywanego brzegu, który nie należy do $\mathbb H^2$.
Znaki $K$ odpowiadają sumom kątów ze wzoru
\eqref{eq:gb-trojkat}.}
\label{fig:gb-trojkaty-modele}
\end{figure}

\begin{przyklad}[Sfera i torus: całka nie zależy od metryki]
\label{prz:gb-sfera-torus}
Na sferze $S_R^2$ stała krzywizna $R^{-2}$ i pole
$4\pi R^2$ dają całkę $4\pi$; promień się skraca.
Na torusie obrotowym niech $a>b>0$, a parametryzacja
ma postać
\[
 X(u,v)=((a+b\cos v)\cos u,
          (a+b\cos v)\sin u,b\sin v),
 \qquad (u,v)\in[0,2\pi]^2.
\]
Różniczkując $X$ znajdujemy
$g=(a+b\cos v)^2\dd u^2+b^2\dd v^2$ i
$\dd A=b(a+b\cos v)\dd u\dd v$.
Dla metryki postaci $\dd s^2+f(s)^2\dd u^2$
rachunek z symboli Christoffela daje
$K=-f''(s)/f(s)$: jedyne niezerowe symbole to
$\Gamma^s_{uu}=-ff'$ oraz
$\Gamma^u_{su}=\Gamma^u_{us}=f'/f$, więc
$g(R(\partial_s,\partial_u)\partial_u,\partial_s)
=-ff''$, a mianownik we wzorze na $K$ wynosi $f^2$.
Biorąc $s=bv$ i $f(s)=a+b\cos(s/b)$, otrzymujemy
\[
 K(u,v)=\frac{\cos v}{b(a+b\cos v)},\qquad
 \int_{T^2}K\dd A
 =\int_0^{2\pi}\!\int_0^{2\pi}\cos v\,\dd u\dd v=0.
\]
Krzywizna dodatnia po zewnętrznej i ujemna po
wewnętrznej stronie torusa znoszą się w całce.
Zgadza się to z $\chi(T^2)=0$. Zmiana metryki
zmienia lokalne $K$ i $\dd A$, ale nie może zmienić
ich całki na tej samej zwartej powierzchni.
\end{przyklad}

\subsection*{Uogólnienie I: krzywizna skupiona w punktach}

Na gładkiej powierzchni krzywizna $K\,\dd A$ jest
rozłożona na obszarze. Można jednak skleić płaskie
wielokąty w zamkniętą powierzchnię: wnętrze każdej ściany
i punkt wewnętrzny każdej sklejonej krawędzi ma płaskie
otoczenie. Cała krzywizna jest wtedy skupiona w
\emph{wierzchołkach}. Przy wierzchołku $v$ sumujemy kąty
przylegających ścian i oznaczamy tę sumę przez $\theta_v$.
Brakujący względem pełnego obrotu kąt
\begin{equation}\label{eq:gb-deficyt-kata}
 \delta_v=2\pi-\theta_v
\end{equation}
nazywamy \emph{deficytem kąta}. Może być ujemny, jeśli
wokół wierzchołka mieści się więcej niż $2\pi$.
Lokalnie jest to powierzchnia stożka o kącie całkowitym
$\theta_v$; zwykły płaski punkt ma $\theta_v=2\pi$.

\begin{twierdzenie}[Wielościenna postać Gaussa--Bonneta]
\label{thm:gb-wielosciany}
Niech $P$ będzie zwartą powierzchnią bez brzegu,
powstałą przez sklejenie skończonej liczby trójkątów
euklidesowych wzdłuż całych boków. Oznaczmy zbiór
jej wierzchołków przez $\mathrm{Vert}(P)$. Wówczas
\begin{equation}\label{eq:gb-wielosciany}
 \boxed{\displaystyle\sum_{v\in\mathrm{Vert}(P)}(2\pi-\theta_v)
                      =2\pi\chi(P).}
\end{equation}
W sumie można pominąć wierzchołki o $\theta_v=2\pi$.
\end{twierdzenie}
\begin{proof}
Niech $V,E,F$ oznaczają liczbę wierzchołków, krawędzi
i trójkątów. Przy każdym $v$ odejmujemy od $2\pi$
wszystkie kąty trójkątów stykających się w $v$.
Każdy kąt pojawia się dokładnie raz, a suma trzech kątów
jednego płaskiego trójkąta wynosi $\pi$. Dlatego
\[
 \sum_v(2\pi-\theta_v)=2\pi V-\pi F.
\]
Każdy trójkąt ma trzy boki, a każda krawędź należy do
dwóch trójkątów, zatem $3F=2E$. W konsekwencji
$2\pi V-\pi F=2\pi(V-E+F)=2\pi\chi(P)$.
Ten rachunek wyjaśnia, dlaczego liczymy kąty \emph{wewnątrz}
ścian, a nie kąty dwuścienne bryły w $\R^3$.
\end{proof}

\begin{przyklad}[Czworościan i sześcian]
W każdym wierzchołku foremnego czworościanu spotykają
się trzy trójkąty równoboczne. Zatem $\theta_v=\pi$,
$\delta_v=\pi$, a cztery wierzchołki dają $4\pi$.
W sześcianie spotykają się trzy kąty proste:
$\theta_v=3\pi/2$ i $\delta_v=\pi/2$; osiem
wierzchołków znowu daje $4\pi$. Obie powierzchnie są
topologicznie sferami, więc $2\pi\chi(S^2)=4\pi$.
Rozkład krzywizny jest inny, lecz suma pozostaje taka sama.
\end{przyklad}

Możemy połączyć krzywiznę gładką z punktowymi deficytami.
Załóżmy, że zwarta powierzchnia $S$ bez brzegu ma skończenie wiele
punktów $p_i$ o otoczeniach będących płaskimi stożkami
o kątach $\theta_i$, a poza nimi jest gładką powierzchnią
riemannowską. Wtedy
\begin{equation}\label{eq:gb-stozki}
 \boxed{\displaystyle
 \int_{S\setminus\{p_1,\ldots,p_N\}}K\,\dd A
       +\sum_{i=1}^N(2\pi-\theta_i)
       =2\pi\chi(S).}
\end{equation}
Istotnie, wytnijmy mały dysk stożkowy przy każdym $p_i$.
Otrzymana gładka powierzchnia ma charakterystykę Eulera
$\chi(S)-N$. Brzeg usuniętego dysku ma długość
$\theta_i r$ i krzywiznę geodezyjną $-1/r$:
minus bierze się z orientacji brzegu \emph{dziury}.
Jego wkład do całki brzegowej to zatem $-\theta_i$.
Podstawienie do~\eqref{eq:gauss-bonnet-brzeg} daje
$\int K\,\dd A-\sum_i\theta_i=2\pi(\chi(S)-N)$,
co po przeniesieniu wyrazów jest~\eqref{eq:gb-stozki}.
Dla powierzchni nieorientowalnej stosujemy ten rachunek
na dwukrotnym nakryciu orientacyjnym i dzielimy przez $2$:
całka, liczba deficytów i charakterystyka Eulera podwajają się.
Gdy część gładka jest płaska, pozostaje dokładnie
\eqref{eq:gb-wielosciany}. W języku miar można myśleć
o krzywiźnie jako o $K\,\dd A+\sum_i\delta_i\delta_{p_i}$:
symbol $\delta_{p_i}$ oznacza masę jednostkową w punkcie,
a współczynnik $\delta_i$ jest deficytem kąta.

\subsection*{Uogólnienie II: parzysty wymiar i twierdzenie Cherna}

W wymiarze dwa $K$ jest całą informacją o krzywiźnie
riemannowskiej w danym punkcie. W wymiarze cztery lub
większym nie wystarczy scałkować samej krzywizny skalarnej:
liczy się określona \emph{kombinacja} składowych tensora
Riemanna. Oto jej precyzyjny opis w języku form.

Niech $(M^{2m},g)$, $m\geq1$, będzie zorientowaną rozmaitością
riemannowską, a $(e_1,\ldots,e_{2m})$ dodatnio zorientowaną
lokalną ramą ortonormalną. Macierz form krzywizny ma wyrazy
\[
 \Omega_{ij}(X,Y)=g(R(X,Y)e_j,e_i),\qquad
 \Omega_{ji}=-\Omega_{ij}.
\]
\emph{Pfaffian} takiej macierzy jest formą stopnia $2m$:
\begin{equation}\label{eq:gb-pfaff-def}
 \operatorname{Pf}(\Omega)=\frac1{2^m m!}
 \sum_{\sigma\in S_{2m}}\operatorname{sgn}(\sigma)
 \Omega_{\sigma(1)\sigma(2)}\wedge\cdots\wedge
 \Omega_{\sigma(2m-1)\sigma(2m)}.
\end{equation}
Zasada konstrukcji: parujemy wszystkie kierunki styczne,
bierzemy iloczyny zewnętrzne odpowiadających im krzywizn
i dodajemy je ze znakami wyznaczonymi przez permutacje.
Dlaczego otrzymujemy globalną formę? Przy zmianie dodatniej
ramki macierzą $Q\in\mathrm{SO}(2m)$ mamy
$\widetilde\Omega=Q^{\mathsf T}\Omega Q$ oraz
$\operatorname{Pf}(Q^{\mathsf T}\Omega Q)
=\det(Q)\operatorname{Pf}(\Omega)=\operatorname{Pf}(\Omega)$.
Tożsamość wynika z rozpisania wzoru permutacyjnego:
po antysymetryzacji iloczyn $2m$ współczynników $Q$
daje jego wyznacznik. Nie ma pochodnych $Q$, bo
$\Omega$ pochodzi z tensora krzywizny. Lokalne formy
sklejają się więc na przecięciach map.
W orientacji przeciwnej znak formy się zmienia razem
ze znakiem całkowania. Dla $m=1$ dostajemy po prostu
$\operatorname{Pf}(\Omega)=\Omega_{12}=K\,\dd A$;
normalizacja z poniższego twierdzenia daje zatem
dokładnie~\eqref{eq:gb-zamknieta}.

\begin{twierdzenie}[Cherna--Gaussa--Bonneta]
\label{thm:chern-gauss-bonnet}
Jeżeli $(M^{2m},g)$ jest zwartą, zorientowaną
rozmaitością riemannowską \emph{bez brzegu}, to
\begin{equation}\label{eq:chern-gauss-bonnet}
 \boxed{\displaystyle
 \chi(M)=\frac1{(2\pi)^m}\int_M
                       \operatorname{Pf}(\Omega).}
\end{equation}
Forma $\operatorname{Pf}(\Omega)/(2\pi)^m$ reprezentuje
w kohomologiach de Rhama rzeczywistą \emph{klasę Eulera}
wiązki stycznej $TM$. Całka nie zależy więc od wyboru
metryki, choć sam Pfaffian zależy od niej punktowo.
\end{twierdzenie}

Dowód przez teorię Cherna--Weila przedstawiamy w
Rozdziale~\ref{ch:chern-weil}. Warto rozdzielić jego kroki:
zamkniętość formy i niezależność jej klasy od koneksji
nie wystarczają jeszcze do rozpoznania klasy Eulera.
Tę identyfikację daje Stwierdzenie~\ref{prop:euler-identification-26},
korzystające z klasy Thoma z Twierdzenia~\ref{thm:euler-thom-16}.
Lemat~\ref{lem:angular-form-26} ustala normalizację formy kątowej:
jej całka po każdym włóknie sferycznym wynosi $1$.
Twierdzenie Stokesa wiąże następnie całkę z Pfaffianu
z sumą indeksów zer pola stycznego, równą $\chi(M)$.
Tu wykonujemy dwa testy wzoru: przypadek $m=1$ powyżej
oraz następujący przykład czterowymiarowy.

W wymiarze cztery
\begin{equation}\label{eq:gb-pfaff-4}
 \operatorname{Pf}(\Omega)=
 \Omega_{12}\wedge\Omega_{34}
 -\Omega_{13}\wedge\Omega_{24}
 +\Omega_{14}\wedge\Omega_{23}.
\end{equation}
Po rozpisaniu w ortonormalnej ramie, przy normach
$|\mathrm{Rm}|^2=\sum_{i,j,k,l}R_{ijkl}^2$ i
$|\mathrm{Ric}|^2=\sum_{i,j}\mathrm{Ric}_{ij}^2$,
otrzymujemy równoważną postać. Wskażmy algebraiczny krok
prowadzący do współczynników $1,-4,1$. Niech $\eta^a$
będzie ramką dualną i $\varepsilon_{1234}=1$ symbolem
całkowicie antysymetrycznym. Ponieważ
$\Omega_{ij}=\frac12\mathcal R_{abji}\eta^a\wedge\eta^b$,
wzór permutacyjny daje
\[
 \operatorname{Pf}(\Omega)
 =\frac1{32}\varepsilon_{ijkl}\varepsilon_{abcd}
       \mathcal R_{abji}\mathcal R_{cdlk}\,\dd\mu_g
 =\frac18\bigl(|\mathrm{Rm}|^2-4|\mathrm{Ric}|^2
                         +\mathrm{Scal}^2\bigr)\,\dd\mu_g.
\]
W ostatnim przejściu rozwijamy
$\varepsilon_{ijkl}\varepsilon_{abcd}
=\det(\delta_{u v})_{u\in(i,j,k,l),\,v\in(a,b,c,d)}$.
Po zebraniu 24 permutacji, użyciu antysymetrii par,
symetrii zamiany par i pierwszej tożsamości Bianchiego
kontrakcja wynosi
$4|\mathrm{Rm}|^2-16|\mathrm{Ric}|^2+4\mathrm{Scal}^2$.
Iloczyn czynnika $1/8$ z $(2\pi)^{-2}$ daje zatem
\begin{equation}\label{eq:gb-4d}
 \chi(M^4)=\frac1{32\pi^2}\int_{M^4}
 \bigl(|\mathrm{Rm}|^2-4|\mathrm{Ric}|^2
                    +\mathrm{Scal}^2\bigr)\,\dd\mu_g.
\end{equation}
Obecność trzech składników pokazuje, czemu
$\int\mathrm{Scal}\,\dd\mu_g$ sama nie może być
czterowymiarowym odpowiednikiem $\int K\,\dd A$.

\begin{przyklad}[Iloczyn dwóch sfer]
Na $M=S_a^2\times S_b^2$ wybierzmy ramę $(e_1,e_2)$
styczną do pierwszej sfery i $(e_3,e_4)$ styczną do
drugiej. W metryce iloczynowej jedyne niezerowe formy
krzywizny występujące w~\eqref{eq:gb-pfaff-4} to
$\Omega_{12}=a^{-2}\dd A_a$ i
$\Omega_{34}=b^{-2}\dd A_b$. Stąd
\[
 \frac1{(2\pi)^2}\int_M\operatorname{Pf}(\Omega)
 =\frac1{4\pi^2}
   \left(\frac{4\pi a^2}{a^2}\right)
   \left(\frac{4\pi b^2}{b^2}\right)=4.
\]
Zgadza się to z $\chi(S^2\times S^2)=4$:
komórki iloczynowego rozkładu CW to jedna $0$-komórka,
dwie $2$-komórki i jedna $4$-komórka. Promienie ponownie
znikają z wyniku.
\end{przyklad}

\subsection*{Brzeg, niezwartość i nieparzysty wymiar}

Dla zwartej, zorientowanej rozmaitości riemannowskiej
parzystego wymiaru o gładkim brzegu twierdzenie Cherna
ma wersję z poprawką brzegową. Jej precyzyjną definicję
i dowód podajemy w
Twierdzeniu~\ref{thm:CGB-boundary-26}; tutaj zapisujemy wzór:
\begin{equation}\label{eq:chern-gb-brzeg-idea}
 \chi(M^{2m})=
 \frac1{(2\pi)^m}\int_M\operatorname{Pf}(\Omega)
 +\int_{\partial M}\mathcal B_g.
\end{equation}
Forma $\mathcal B_g$ na gładkim brzegu zawiera jego
geometrię wewnętrzną i sposób zanurzenia w $M$
(drugą formę fundamentalną). Nie jest w ogólności
samą krzywizną Gaussa brzegu. W wymiarze dwa mamy
$\mathcal B_g=k_g\,\dd s/(2\pi)$; dla brzegu
z narożnikami dochodzą jeszcze skręty opisane
w~\eqref{eq:gauss-bonnet-brzeg}. Zasadniczy test:
domknięta płaska kula $\overline B^{2m}\subset\R^{2m}$
ma zerowy Pfaffian we wnętrzu, lecz $\chi(\overline B^{2m})=1$,
zatem cały wkład w~\eqref{eq:chern-gb-brzeg-idea}
pochodzi z $\partial\overline B^{2m}$.

Bez zwartości całka po wnętrzu również nie musi
odtwarzać $\chi$: dla całej płaskiej $\R^2$
mamy $\int_{\R^2}K\,\dd A=0$ i $\chi(\R^2)=1$.
Na wyczerpujących ją dyskach $B_R$ brakujący wkład
jest widoczny jako $\int_{\partial B_R}k_g\,\dd s=2\pi$.
W twierdzeniach dla przestrzeni niezwartych trzeba
kontrolować granicę takiego wyrazu przy $R\to\infty$;
nie wolno go po prostu skreślić.

Dla porównania, każda zwarta orientowalna rozmaitość
\emph{bez brzegu} nieparzystego wymiaru ma $\chi(M)=0$.
Wynika to z dualności Poincarégo nad $\Q$:
$b_i=b_{n-i}$, a gdy $n$ jest nieparzyste,
składniki $(-1)^i b_i$ i $(-1)^{n-i}b_{n-i}$
w sumie Eulera znoszą się parami. Nie ma więc
w nieparzystym wymiarze analogicznej nietrywialnej
całki z Pfaffianu stopnia najwyższego.

\begin{uwaga}[Co mówi, a czego nie mówi całkowita krzywizna?]
Dla zamkniętej powierzchni orientowalnej rodzaju $g=0$
całka $K$ jest dodatnia, dla $g=1$ zerowa, a dla
$g\geq2$ ujemna. Nie wynika stąd, że $K$ ma ten sam
znak w każdym punkcie: wzór określa sumę wkładów.
Wzór z samym $K$ jest \emph{dwuwymiarowy};
uogólnienie Cherna dotyczy zwartych rozmaitości
riemannowskich parzystego wymiaru. Nie stosujemy
żadnego z tych zapisów bezpośrednio do niezwartych
czterowymiarowych metryk \emph{lorentzowskich}
Schwarzschilda i Kerra: potrzebne są inne założenia,
odpowiednia wersja dla sygnatury metryki i analiza brzegu
bądź zachowania w nieskończoności.
\end{uwaga}

\par\smallskip
{\footnotesize\noindent\emph{Dalsza lektura:}
\href{https://www.wim.uni-mannheim.de/media/Lehrstuehle/wim/schmidt/FSS2024/Riemannian_Geometry/Web/RGch5.html}{krzywizna},
\href{https://davidtong.org/teaching/general-relativity/grhtml/S6}{Schwarzschild i Kerr},
\href{https://link.springer.com/article/10.1140/epjc/s10052-025-13983-8}{kontrakcja Kerra},
\href{https://math.stanford.edu/~ralph/math215b/Milnor.pdf}{Milnor},
\href{https://web.math.utk.edu/~langford/RG.pdf}{porównania},
\href{https://math.mit.edu/~lurie/papers/hadamard.pdf}{Cartan--Hadamard},
\href{https://math.mit.edu/~lguth/Exposition/systmet.pdf}{skalar i małe kule},
\href{https://www.math.uchicago.edu/~may/REU2022/REUPapers/Hale.pdf}{Gauss--Bonnet},
\href{https://arxiv.org/abs/1111.4972}{Cherna--Gaussa--Bonneta},
\href{https://pages.uoregon.edu/gilkey/P251-Gauss-Bonnet-v7a.pdf}{uogólnienie z brzegiem}.\par}
